\documentclass[11pt]{article}
\usepackage[margin=1in]{geometry}
\usepackage{amsmath,amssymb,amsthm,mathtools}
\usepackage[T1]{fontenc}
\usepackage{lmodern}
\usepackage{microtype}
\usepackage{enumitem}
\usepackage{xurl}
\usepackage[colorlinks=true,linkcolor=blue,citecolor=blue,urlcolor=blue]{hyperref}
\hypersetup{pdftitle={Frobenius product, parallel circuits, and sample-complexity barriers for deep Gaussian Bayesian networks},pdfauthor={ChatGPT-6 Astra; Kaarel Haenni; Dmitry Vaintrob}}
\newtheorem{theorem}{Theorem}[section]
\newtheorem{lemma}[theorem]{Lemma}
\newtheorem{proposition}[theorem]{Proposition}
\newtheorem{corollary}[theorem]{Corollary}
\newtheorem{definition}[theorem]{Definition}
\theoremstyle{remark}
\newtheorem{remark}[theorem]{Remark}
\newcommand{\R}{\mathbb R}
\newcommand{\E}{\mathbb E}
\newcommand{\Prb}{\mathbb P}
\newcommand{\Lip}{\operatorname{Lip}}
\newcommand{\sign}{\operatorname{sign}}
\newcommand{\poly}{\operatorname{poly}}
\newcommand{\TC}{\mathsf{TC}}
\newcommand{\NC}{\mathsf{NC}}
\newcommand{\Ppoly}{\mathsf{P}/\mathrm{poly}}
\newcommand{\MAJ}{\operatorname{MAJ}}
\newcommand{\KL}{\operatorname{KL}}
\newcommand{\Gain}{\mathcal G}
\newcommand{\norm}[1]{\left\lVert #1\right\rVert}
\title{Frobenius product, parallel circuits, and sample-complexity\
barriers for deep Gaussian Bayesian networks}
\author{ChatGPT-6 Astra \and Kaarel H\"anni \and Dmitry Vaintrob\thanks{Authors are listed in alphabetical order.}}
\date{Research draft --- 25 October 2026}

\begin{document}
\maketitle
\begin{abstract}
We prove that a Boolean function with a globally unit-margin neural
representation of layer Frobenius product $P$ has a nonuniform threshold
circuit of size polynomial in $d+P$ and depth $O(\log(2+P))$. The original
network may have arbitrary finite depth and widths, and each neuron may
use its own arbitrary zero-preserving $1$-Lipschitz activation. Thus
polynomial product implies nonuniform $\TC^1$, and hence $\NC^2$.
The proof combines a two-ball covering lemma, finite-set sampling of a
short prefix, and bounded-precision lookup tables; it does not require
computable activations or well-conditioned intermediate layers.

For the original tanh Gaussian posterior conditioned only on a global
absolute margin and exact training labels, we prove a sample-to-circuit
converse uniformly at every finite scale $0<\sigma^2\le1/N$.
The circuit resources are polynomial in input dimension, sample count,
and learner depth. The theorem requires logarithmic depth padding
and uses a width threshold exponential in input dimension.
Its proof combines finite canonical features, Gaussian block covers,
and factorially convergent parallel Picard iteration. The block covers
keep the number of representatives polynomial even when the common
prior prefactor approaches one. All confidence-failure tolerances
$\varepsilon<1/2$ are allowed. A separate above-critical product
ceiling gives circuit resources independent of learner depth.
Consequently, for any prescribed prior-scale schedule in the stated
range, a universal sample bound polynomial jointly in teacher
resources and learner depth would imply
$\Ppoly=\TC^1$. Conversely, the stronger assumption
$\Ppoly=\TC^0$ suffices for that joint polynomial learning bound,
by compiling one universal circuit-evaluation family.
An explicit finite-horizon compiler shows that the necessary collapse
already follows for teachers of width $24$ and weights bounded by one,
while the learner width remains unrestricted.
In the opposite direction, previously established
learning bounds imply that any fixed stretched-exponential sample lower bound
$\exp(\Omega(L^\alpha))$, $\alpha>0$,
with teacher resources polynomial in depth $L$ would imply
$\NC\ne\Ppoly$. These are barriers to the two proposed resolutions,
not a resolution of unrestricted-depth learning itself.
\end{abstract}

\section{Introduction}

A deep network can have small individual weights while its product of
layer norms grows exponentially with depth. These two resource measures
lead to different questions. A bound on width, depth, and individual
weights gives a finite description of a globally margined Boolean
classifier. A bound on the product of Frobenius norms imposes a stronger
restriction: we show that it gives a small parallel circuit, regardless
of the network's original depth.

The geometric reason is simple. A linear map of Frobenius norm $M$
sends every Euclidean ball of radius $r$ into the union of two balls of
radius $Mr/\sqrt2$. Coordinatewise $1$-Lipschitz maps preserve this
property. After $O(\log P)$ layers, a network of total product $P$ has
only polynomially many prefix states relevant to its output signs.
The remaining layers can be replaced by a finite classification of
these states. A separate sampling argument bounds the width of the
retained prefix. Normalizing each layer before numerical simulation
removes any dependence on the sizes of intermediate signals.

The result has consequences for Bayesian learning. In the original
Gaussian tanh model, the prior is conditioned on
$|z(x)|>1/2$ at every Boolean input and on the observed labels.
Above Gaussian exponent one, an existing finite-scale posterior ceiling
places most posterior mass on Boolean functions with product polynomial
in the input dimension and sample count, uniformly in learner depth.
Combining that ceiling with our circuit theorem gives a converse: a
distribution-free sample guarantee forces the target to have a small
parallel circuit. Two finite majority constructions handle posterior
mixtures and confidence-failure tolerances arbitrarily close to $1/2$.

The critical exponent one requires a different representation theorem.
We construct finite canonical features preserving every query's
correlations with a common anchor set, without preserving the entire
kernel. Their dimension grows additively with depth. The resulting
recursion is a linear contraction sequence plus residual maps of small
total Lipschitz constant, and parallel Picard iteration evaluates it in
logarithmic threshold depth. Uniform importance quadrature makes each
map a finite circuit despite arbitrary controlled densities and
ill-conditioned anchor spans. A finite-event entropy certificate then
gives an actual posterior ceiling on the full cube. The resulting
converse includes $L$ in its polynomial resources, assumes
$L\ge C\log(d+n+L+2)$, and chooses width above
$2^{Cd}\poly(d+n,L)$. These qualifications still permit the universal
polynomial-bound implication below.

Fixed prefactors $0<t<1$ need one more argument, since their exponent
gaps shrink with width. For the two largest eigenvalues of a
controlled row's second moment, Gaussian entropy bounds its
two-ball covering radius more strongly than its ordinary Lipschitz
gain. The strict gap accumulates over a short prefix while charging
the entire suffix. A normalized importance rule evaluates that
prefix without representing large or small physical scales.
This first gives a converse uniform on each range $0<t\le t_0<1$.
To cover the moving boundary $t\to1$, we group
$O((t^{-2}-1)^{-1})$ layers at a time. On most blocks the linear
map is a contracting isometry plus a Hilbert--Schmidt perturbation
of bounded size, and the nonlinear residual is small. Each such
block admits a constant-size cover. Only logarithmically many
blocks are needed, while parallel Picard iteration evaluates their
potentially long prefix. The final theorem is uniform over every
$0<t\le1$, equivalently every $0<\sigma^2\le1/N$.

There are two separate consequences. A universal joint polynomial
sample upper bound would collapse $\Ppoly$ to nonuniform $\TC^1$.
The stronger collapse $\Ppoly=\TC^0$ is sufficient for that
learning bound: one universal evaluator fixes the resource
exponent for all bounded-weight teachers.
Any fixed stretched-exponential depth lower bound, when both input dimension and
bounded-weight teacher resources are polynomial in depth, would
separate $\Ppoly$ from every nonuniform polylogarithmic-depth Boolean
circuit class. Neither consequence is currently an unconditional
complexity theorem. The arguments also leave room for intermediate
sample growth between polynomial and exponential.
Appendix~\ref{app:constant-width} shows that the necessary collapse
already applies to fixed-width teachers: a width-$24$ tanh network with coefficients bounded by one can simulate any small Boolean circuit at polynomially
larger depth. Its memory is exponentially delicate, so this observation
does not supply a robust or polynomial-product representation.

The geometric, canonical-feature, finite-event extraction, and
circuit reductions are proved below. We state the previously established
inputs explicitly: the finite-scale posterior ceiling in
Section~\ref{sec:ceiling}, the scalar Bellman facts and finite-width
Fourier target mass in Section~\ref{sec:critical-parallel}, the companion
scale-change target mass in Section~\ref{sec:fixed-prefactor}, and the
circuit-target learning bound in Section~\ref{sec:negative}. The
critical endpoint is proved separately, rather than inferred by taking
a limit in the above-one result.
The conditional sufficiency direction also uses the companion
normalized intrinsic learning theorem, stated explicitly where needed.

\paragraph{Relation to previous work.}
Norm-controlled generalization and compression are well established.
Golowich, Rakhlin, and Shamir~\cite{grs} bound Rademacher complexity
under layer norm constraints, improving depth dependence and obtaining
size-independent bounds under additional assumptions.
Their full-version Theorem~3 already obtains a rank-one layer
replacement and a Lipschitz suffix; its displayed uniform error
decreases polynomially with the retained prefix length.
Here the global Boolean margin is preserved using a logarithmic
prefix and polynomially many covering states.
Sellke~\cite{sellke} refines the Frobenius-product analysis for ReLU
networks, often removing explicit depth dependence.
Our sampling and normalization arguments belong to this broader
norm-compression framework. The conclusion studied here is an exact
nonuniform circuit for a globally margined Boolean classifier, followed
by a comparison with the original conditioned Bayesian posterior.
A norm-based generalization estimate alone does not make that
posterior-to-teacher-complexity comparison.

Circuit-complexity barriers to neural depth separation also precede
this work. Vardi and Shamir~\cite{vardi-shamir} relate deeper ReLU
approximation lower bounds to circuit lower bounds and natural-proof
barriers, using a reduction to polynomially bounded weights.
Vardi, Reichman, Pitassi, and Shamir~\cite{vardi-benign} establish
related barriers for approximating efficiently computable functions
with polynomial Lipschitz bounds. The implications here concern a
different quantitative object: distribution-universal posterior
sample complexity for bounded-resource tanh Boolean teachers.
We make no claim that connecting neural representation bounds to
circuit complexity, or compressing norm-bounded networks, is new.

Jacot~\cite{jacot-rank,jacot-bottleneck} studies a related depth
baseline for homogeneous networks: the leading representation
cost is governed by a nonlinear rank, with finite-depth corrections
and quantitative bounds on deviation from low-rank layer matrices.
Those results help distinguish low-dimensional representation from
the complexity of its encoder and decoder. They do not by themselves
bound that residual complexity by a polynomial in the teacher
resources considered here.

\section{Networks and circuit conventions}
\label{sec:conventions}

Let $X_d=\{-1,1\}^d$ and $f:X_d\to\{-1,1\}$. A network of weighted
depth $D\ge1$ has augmented input $h^0(x)=(1,x)$, hidden layers
\[
 h^s_i(x)=\phi_{s,i}\bigl((W^s h^{s-1}(x))_i\bigr),
 \qquad 1\le s<D,
\]
and scalar linear output $z(x)=W^D h^{D-1}(x)$. Every hidden width is
finite. Unless specified otherwise, each $\phi_{s,i}:\R\to\R$ satisfies
\[
 \phi_{s,i}(0)=0,\qquad
 |\phi_{s,i}(u)-\phi_{s,i}(v)|\le |u-v|.
\]
There are no additional affine biases: constants enter through the
augmented input and may be propagated by ordinary hidden units.
Write
\[
 P_F(W)=\prod_{s=1}^D\norm{W^s}_F.
\]
A global unit-margin representation means $f(x)z(x)\ge1$ for every
$x\in X_d$. We allow arbitrary real weights and arbitrary activations
in this definition; no computability assumption is made.

All circuit classes in this paper are \emph{nonuniform}. Boolean
circuits use a fixed bounded-fan-in basis, for example
$\{\mathsf{AND},\mathsf{OR},\mathsf{NOT}\}$, together with constants.
Threshold circuits may equivalently use unbounded-fan-in majority
gates and bounded-fan-in Boolean gates. Our constructed weighted
threshold gates have polynomially bounded integer weights and can be
converted to majority gates with polynomial overhead. Thus $\TC^1$
denotes polynomial size and logarithmic depth in the ordinary
nonuniform sense. We use
\[
 \NC=\bigcup_{k\ge1}\NC^k,
\]
where $\NC^k$ has polynomial size and depth $O(\log^k d)$.
Size may count wires and input nodes; retaining all $d$ named input
nodes changes our bounds only by an additive $d$.

\begin{theorem}[Product-to-parallel-circuit theorem]
\label{thm:main}
There are universal constants $a,c>0$ such that the following holds.
If $f:X_d\to\{-1,1\}$ has a global unit-margin representation as above
with $P_F(W)\le P$, then $P\ge1$ and $f$ has a threshold circuit with
at most $(d+P+2)^a$ wires and depth at most $c\log(P+2)$.
It also has a bounded-fan-in Boolean circuit of polynomial size in
$d+P$ and depth at most $c\log^2(P+2)$, after changing $a,c$.

The constants are independent of the original depth, widths, physical
weights, and choices of activations. In particular, any family with
$P=\poly(d)$ belongs to nonuniform $\TC^1\subseteq\NC^2$.
\end{theorem}

The theorem is about represented functions. It does not give an
efficient procedure for finding a circuit from a network, and does not
put the evaluation problem for supplied arbitrary real networks in
$\NC$. Its proof deliberately permits auxiliary activations outside
the fixed-tanh class.

\section{A short prefix suffices}

\subsection{Relevant inputs}

\begin{lemma}[Relevant-coordinate reduction]
\label{lem:junta}
Every globally unit-margin representation has Frobenius product at least one.
Its Boolean function depends on at most $P_F(W)^2$ input coordinates.
After fixing the other coordinates, it has a representation on just
the relevant coordinates whose product is no larger.
\end{lemma}
\begin{proof}
Zero preservation and the Lipschitz bounds imply
\[
 \E_{x\sim\mathrm{Unif}(X_d)}z(x)^2
 \le \left(\prod_{s=2}^D\norm{W^s}_F^2\right)
       \E\norm{W^1(1,x)}_2^2
 =P_F(W)^2.
\]
Here $\E[(1,x)(1,x)^T]=I$, and the formula also covers $D=1$.
The margin makes the left side at least one.

Put $T=\prod_{s=2}^D\norm{W^s}_F$, with an empty product equal to one.
For each relevant coordinate $i$, some pair of inputs differing only
at $i$ has opposite labels. Their outputs differ by at least two, so
\[
 2\le 2T\norm{W^1 e_i}_2.
\]
Summing the resulting squared inequalities over relevant coordinates
gives their number $k\le P_F(W)^2$.

For irrelevant coordinates choose independent uniform signs $a_i$.
Absorb them into the first bias column by replacing $w_0$ with
$w_0+\sum_i a_i w_i$, and delete their columns. The expected squared
Frobenius norm of the new first matrix equals the old squared norm.
Some choice therefore does not increase it. Irrelevance ensures that
all these restrictions preserve the Boolean function and its margin.
\end{proof}

\subsection{The two-ball lemma}

\begin{lemma}[Dimension-free two-ball covering]
\label{lem:two-ball}
For every finite real matrix $W$, center $a$, and $r\ge0$, and every
coordinatewise $1$-Lipschitz map $\Phi$, the set
$\Phi(W B(a,r))$ is covered by at most two Euclidean balls of radius
$r\norm W_F/\sqrt2$.
\end{lemma}
\begin{proof}
Zero radius and the zero matrix are immediate. Translation and scaling
reduce the linear part to $W B(0,1)$ with $\norm W_F=1$.
Let $\alpha$ be its largest singular value. If $\alpha^2\le1/2$,
the operator-norm bound gives one ball. Otherwise put
\[
 b^2=1-\alpha^2,\qquad
 c=\frac{\alpha^2-b^2}{2\alpha}
   =\frac{2\alpha^2-1}{2\alpha},
\]
and let $u$ be a leading left singular vector. The remaining singular
part has operator norm at most $b$. For a unit-ball input whose leading
right-singular coordinate is $t$, distance squared from its image to
the center $cu$ or $-cu$ matching the sign of $t$ is at most
\[
  (\alpha |t|-c)^2+b^2(1-t^2).
\]
This is convex as a function of $|t|\in[0,1]$. Both endpoint values are
\[
 c^2+b^2=(\alpha-c)^2=\frac1{4\alpha^2}\le\frac12.
\]
The two centers therefore give the claimed linear cover. Apply
$\Phi$ to the translated and rescaled centers; its Euclidean
Lipschitz constant is at most one. The centers need not belong to the
image, which does not affect the covering assertion.
\end{proof}

\begin{lemma}[Margin-relevant prefix states]
\label{lem:prefix}
Let a network of product $P$ have global unit margin on a set of inputs
of Euclidean norm at most $R$. Put
$S_j=\prod_{s\le j}\norm{W^s}_F$. For $j\le D-1$, its $j$th hidden
image has a cover by at most $2^j$ balls of radius
\[
 r_j=R S_j 2^{-j/2}.
\]
If
\[
 j=\left\lceil2\log_2(16PR)\right\rceil\le D-1,
\]
one can select at most $512P^2R^2$ actual normalized prefix states
$c_a=h^j(x_a)/S_j$, each with its label $f(x_a)$, so that every
$u=h^j(x)/S_j$ has a same-label selected state within $1/(8P)$,
while every selected state of the opposite label is at distance at
least $2/P$. All these states have norm at most $R$.
\end{lemma}
\begin{proof}
Iterate Lemma~\ref{lem:two-ball}, starting from the radius-$R$ input
ball. The remaining suffix is Lipschitz with constant at most $P/S_j$.
Two hidden images in one covering ball therefore have output
difference at most
\[
 2r_jP/S_j=2PR2^{-j/2}\le1/8,
\]
whereas opposite labels require difference at least two. Each nonempty
ball is thus monochromatic. Select an actual image in each nonempty
ball; this doubles the cover radius. Divide by $S_j$ to get the
same-label bound. The suffix Lipschitz bound gives the opposite-label
bound directly. Finally $2^j\le2(16PR)^2$, and zero preservation gives
$\norm{h^j(x)}\le S_jR$.
\end{proof}

\subsection{Sampling only the retained prefix}

We next bound the width of the short prefix. The following argument
uses the freedom to change the neuronwise activations; it is not a
fixed-activation tanh compression claim.

\begin{lemma}[Finite-set prefix compression]
\label{lem:maurey}
Let $X$ be a nonempty finite set of inputs of norm at most $R$, and let a network
as in Section~\ref{sec:conventions} have product at most $P$.
For $1\le h\le D-1$ and $\epsilon>0$, its first $h$ hidden layers can
be replaced by layers of width at most
\[
 m=\left\lceil
 \max\{1,64(PRh/\epsilon)^2(1+\log|X|)\}
 \right\rceil
\]
so that all final outputs on $X$ change by at most $\epsilon$ and the
full product never increases. The replacement activations remain
zero-preserving and $1$-Lipschitz. No bound depends on the number or
widths of the later layers.
\end{lemma}
\begin{proof}
Consider consecutive matrices $W,V$. Let $w_i$ be row $i$ of $W$ and
$v_i$ column $i$ of $V$, and put $a_i=\norm{w_i}_2$ and
$b_i=\norm{v_i}_2$. Terms with $a_i b_i=0$ contribute zero and may be
deleted. If all contributions are zero, the intervening computation
can be replaced by a zero layer. Otherwise set
\[
 A=\sum_i a_i b_i\le\norm W_F\norm V_F=:B,
 \qquad p_i=\frac{a_i b_i}{A},\qquad u_i=\frac{w_i}{a_i},
 \qquad \psi_i(t)=\frac{\phi_i(a_i t)}{a_i}.
\]
The activation $\psi_i$ has the required properties. The preactivation
of $V$ is $\sum_i a_i v_i\psi_i(u_i^T h)$.
Sample $m$ indices independently with probabilities $p_i$. For each
sampled index use incoming row $u_i$, activation $\psi_i$, and outgoing
column $A v_i/(m b_i)$. The new preactivation is an unbiased empirical
average of the random vectors
\[
 Y(h)=\frac{A v_i}{b_i}\psi_i(u_i^T h).
\]
For every draw, the two new Frobenius norms are $\sqrt m$ and
$A/\sqrt m$, so their product is $A\le B$.

Let $S$ be the product of all preceding layer norms. On $X$ one has
$\norm h\le SR$, hence $\norm{Y(h)}\le ASR=:J$.
The finite-input Hilbert sampling estimate gives a common sample
with every preactivation error smaller than $\delta$ whenever
\[
 64J^2(1+\log|X|)\le m\delta^2,\qquad m\ge1.
\]
Here is a finite-sum proof with deliberately loose constants. If
$F_x$ is the norm of the centered sum of $m$ sampled vectors, the
Hilbert variance identity gives $\E F_x\le J\sqrt m$.
Changing one index changes $F_x$ by at most $2J$. A finite
bounded-differences argument, obtained by iterating the scalar
bound $\E e^{t(Z-\E Z)}\le e^{t^2c^2/2}$ for a variable of
oscillation at most $c$, gives
\[
 \E e^{t(F_x-\E F_x)}\le e^{2t^2mJ^2}.
\]
The scalar bound follows from convexity of the exponential and
$\cosh u\le e^{u^2/2}$. Summing these exponential potentials
over $X$ and taking $t=\delta/(8J^2)$ proves the displayed
sampling estimate: its budget bounds the mean by $m\delta/2$
and makes the total exponential potential strictly smaller
than $e^{tm\delta/2}$. Thus some common sample has
$F_x<m\delta$ for every $x$. This argument includes $|X|=1$.
If $J=0$, all sampled features vanish. A zero full product
makes the final network output identically zero. Both cases
are handled directly.

When the current product and $SB$ are positive, the suffix after
$V$, including its activation if present, is Lipschitz with
constant at most $P/(SB)$. Take
$\delta=(\epsilon/h)SB/P$. Since $A\le B$, the displayed
choice of $m$ guarantees that the final output error in this
replacement is at most $\epsilon/h$. Zero-product networks
are replaced by zero networks before making any division.
Process the first $h$ hidden layers successively, allotting error
$\epsilon/h$ to each. The product does not increase at any step, and
altering the next matrix never changes an already bounded hidden
width. Summing the output errors proves the assertion.
\end{proof}

\section{Compiling the prefix at bounded precision}
\label{sec:compiler}

\begin{proof}[Proof of Theorem~\ref{thm:main}]
Apply Lemma~\ref{lem:junta}. If $f$ is constant, a constant gate
suffices. Otherwise retain $k\le P^2$ relevant inputs and put
$R=\sqrt{k+1}$. Let $\overline P=2P$ and
\[
 j_0=\left\lceil2\log_2(16\overline P R)\right\rceil,
 \qquad h=\min\{j_0,D-1\}.
\]
Apply Lemma~\ref{lem:maurey} with error $1/4$ to these $h$ layers
(nothing is required if $h=0$), then double the final readout. The new
network has signed margin at least $3/2$ and product at most
$\overline P$. Its retained widths are at most a universal polynomial
$C=C(P)$, since $R\le\sqrt{P^2+1}$,
$1+\log|X_k|=1+k\log2$, and $h=O(\log(P+2))$.

If $D-1\ge j_0$, apply Lemma~\ref{lem:prefix} to the new network
using the upper bound $\overline P$. It supplies polynomially many
normalized centers with the stated separation. If $D-1<j_0$, we will
evaluate the entire short network and its linear readout instead.
In both cases only $h\le j_0$ hidden layers need be simulated.

Write $M_s=\norm{W^s}_F$, $S_s=\prod_{i\le s}M_i$, and
$A_s=W^s/M_s$ for the new network. All $M_s$ are nonzero because its
margin is nonzero. Define $u_s=h^s/S_s$ and, neuron by neuron,
\[
 \Psi_{s,i}(v)=\frac{\phi_{s,i}(S_s v)}{S_s}.
\]
Then the exact normalized computation is
\[
 u_s=\Psi_s(A_su_{s-1}),\qquad
 \norm{A_s}_F=1,\qquad \norm{u_s}_2\le R.
\]
Every $\Psi_{s,i}$ is zero-preserving and $1$-Lipschitz, even when
$S_s$ is arbitrarily large or small.

Put $m_0=\max\{C,k+1,2\}$. Choose positive dyadic $\eta,\delta$,
each within a factor two of and no larger than the corresponding value
\[
 \frac1{4096\overline P(R+1)(j_0+1)},\qquad
 \frac1{4096\overline P m_0(j_0+1)}.
\]
Round each $A_s$ entry on a dyadic mesh between
$\eta/(2m_0)$ and $\eta/m_0$, so its Frobenius error is at most
$\eta$. Set $B_0=\lceil R+2\rceil$. Replace each $\Psi_{s,i}$ on
$[-B_0,B_0]$ by a lookup table with grid spacing $\delta$ and table
values rounded to multiples of $\delta$. Lipschitzness gives scalar
error at most $2\delta$. Its table length, all inverse meshes, and all
integer ranges are polynomial in $P+2$; every stored number has
$O(\log(P+2))$ bits. This observation uses only existence of the finite
table, not computability of the activation.

If $e_s$ is the Euclidean error of the simulated state, then
\[
 e_s\le(1+\eta)e_{s-1}+\eta R+2\sqrt C\,\delta,
 \qquad e_0=0.
\]
The displayed choices give $e_s<1/(64\overline P)$ for every
$s\le h$: for example, $(1+\eta)^{j_0}<2$ and the accumulated
right-hand error is at most
$6/[4096\overline P]$. Rounded preactivations stay inside
$[-B_0,B_0]$, as required.

In the long-network case, round each normalized center to Euclidean
error at most $1/(64\overline P)$, using a comparable dyadic mesh
in each coordinate. A correct-label center then has computed distance
at most $5/(32\overline P)$, whereas every opposite-label center has
computed distance at least $63/(32\overline P)$. Choose a dyadic number
\[
 \tau\in[1/(2\overline P^2),1/\overline P^2].
\]
Compute the squared distances of the rounded vectors exactly, and
output $+1$ precisely when some positive-labeled center has squared
distance at most $\tau$. A positive input has such a center, and a
negative input has none. The gap proves exact classification.
In the short-network case, the normalized final readout has norm one
and signed margin at least $1/\overline P$. Round its coefficients
to the same Frobenius precision and take its sign. Its additional
error is at most $\eta R+(1+\eta)e_h<1/(2\overline P)$, so its
sign is also exact.

It remains to justify the circuit bounds without an unproved
transcendental-arithmetic assumption. All table values and rounded
states have a common dyadic grid of inverse-polynomial spacing.
All rounded matrix entries also have polynomially bounded dyadic
denominators. After clearing denominators within one layer, an affine
preactivation is a linear function of the preceding state bits with
polynomially bounded integer weights and polynomial integer range.
Compare this linear function to every table grid boundary by one
threshold layer. Boolean gates identify its grid interval, and, for
each output bit, an OR over those intervals whose stored table value
has that bit equal to one implements the lookup. This has constant
threshold depth and polynomial size. The lookup resets the output to
its fixed grid, so denominators do not accumulate through the layers.
Signed values can be encoded by an offset integer; this only changes
the affine constants.

For distance tests, a difference of two rounded coordinates still
has $O(\log(P+2))$ bits. Its square is a polynomial-size table.
A threshold comparison of the sum of these squares implements each
test, and an OR combines the positive tests. The short-case readout
is itself a threshold gate. Every threshold used here has integer
weights and a threshold of polynomial magnitude. Negative weights
can be moved to complemented inputs and the constant threshold
adjusted. Repeating inputs according to the positive integer weights,
and adding constant inputs, converts such a gate to a majority gate
with polynomial overhead. There are $O(\log(P+2))$ simulated layers,
proving the threshold-circuit assertion.

For completeness, majority on $q$ bits has a bounded-fan-in Boolean
circuit of polynomial size and depth $O(\log(q+2))$. Use constant-depth
full adders to compress three summands to two, in parallel at each bit
position. After $O(\log(q+2))$ compression stages, two binary integers
remain; add and compare them in $O(\log(q+2))$ further depth, even by
a ripple circuit on their $O(\log(q+2))$ bits. Replacing each majority
gate of our threshold circuit this way gives total depth
$O(\log^2(P+2))$ and polynomial size. Finally, reconnect the retained
input coordinates to their original names. This proves the theorem.
\end{proof}

\begin{remark}[What the representation theorem does not assert]
The proof bounds circuit resources, not the original network's
parameters. It does not assert that a small bounded-weight deep
network has polynomial product, or that a majority of low-product
functions has low product. No closure of product norm under majority
will be needed below. The arbitrary-activation compression is used
only as an auxiliary step in compiling the represented Boolean
function; the statistical model remains exactly tanh.
\end{remark}

\begin{corollary}[Other Lipschitz constants]
If the activations in hidden layer $s$ are zero-preserving and have
Lipschitz constants at most $\lambda_s>0$, the conclusion of
Theorem~\ref{thm:main} holds with the effective product
\[
 P_{\mathrm{eff}}=P_F(W)\prod_{s=1}^{D-1}\lambda_s
\]
in place of $P$. In particular, the result includes ReLU.
\end{corollary}
\begin{proof}
Write $\phi_{s,i}=\lambda_s\psi_{s,i}$ with $\psi_{s,i}$
$1$-Lipschitz. Replace the hidden state by $h^s/\lambda_s$, multiply
the following matrix by $\lambda_s$, and use $\psi_{s,i}$ as the
activation. The resulting representation has the same output and
product $P_{\mathrm{eff}}$. Apply Theorem~\ref{thm:main}.
\end{proof}

\section{The original Gaussian posterior and its ceiling}
\label{sec:ceiling}

Fix input dimension $d\ge1$, weighted depth $L\ge3$, and common hidden
width $N$. In this section all hidden activations are tanh. The raw
Gaussian prior has independent entries with variances
\[
 W^1_{ij}:\frac{t^2}{d},\qquad
 W^s_{ij}:\frac{t^2}{N}\ (2\le s<L),\qquad
 v_i:\frac{t^2}{N^2},
\]
where $v=W^L$ is the linear readout. Equivalently, put
$\sigma^2=t^2/N$ and use energy coordinates with squared norm
\[
 \kappa_N(W)=\frac dN\norm{W^1}_F^2+
       \sum_{s=2}^{L-1}\norm{W^s}_F^2+N\norm v_2^2.
\]
Their coordinates have common variance $\sigma^2$.
The exponent convention $\sigma^2=N^{-\gamma}$ corresponds to
$t^2=N^{1-\gamma}$.

For a consistently labeled sample $S=((x_i,y_i))_{i=1}^m$, let $Q_S$
be the raw prior conditioned on
\[
 |z(x)|>1/2\quad(x\in X_d),\qquad
 y_i z(x_i)>0\quad(1\le i\le m).
\]
There are no further constraints on hidden activations, internal
margins, operator norms, or parameter size. Let
\[
 p_S(f(x)\mid x)=Q_S\{f(x)z(x)>0\}.
\]
The confidence-failure mass under an input law $\mathcal D$ is
$\mathcal D\{x:p_S(f(x)\mid x)\le2/3\}$.

We quote the following established result with its explicit constants.
It is the finite-scale posterior energy ceiling in
\cite[section ``A posterior energy ceiling and a necessary
representation condition'']{unrestricted}; the exact source file is
identified in the bibliography.

\begin{theorem}[Imported finite-scale posterior ceiling]
\label{thm:ceiling}
Fix $d,n\ge1$, $L\ge3$, $\zeta>0$, and $0<\beta<1$. Define
\begin{align*}
 C_0&=n+1,\\
 R_0&=\max\{4(n+1),(d+1)\log(16(n+1))\},\\
 M&=\left\lceil\max\{4,C_0,2R_0,dC_0(d+1)R_0^2\}\right\rceil,
 &T&=MC_0R_0^2,\\
 K&=2+(L+14)(2T)^{1/(L-1)},\\
 b^2&=\max\{2(d+1)\log2+4\log8,\ 16\log9+8\log(8L)\},\\
 A&=\left(\frac{K+2}{L-1}\right)^{(L-1)/2},\\
 \xi&=\min\left\{1,\frac1{4\sqrt{L(K+1)}},
                         \frac1{8\sqrt2 LA}\right\},\\
 N_0&=\left\lceil\max\left\{M,2,(b/\xi)^{2/\zeta},
                                      2\log(2/\beta)\right\}\right\rceil.
\end{align*}
For every integer $N\ge N_0$, every actual finite scale
$0<\sigma^2\le N^{-(1+\zeta)}$, and every consistently labeled sample
of at most $n$ inputs, $Q_S$ assigns probability at least $1-\beta$ to
Boolean functions that admit a global unit-margin tanh representation
of product at most $P_0$, where
\[
 \boxed{\quad P_0^2=d^{-1}e^{19}(2T)^{3/2}.\quad}
\]
Such a representation can be taken at the same width and depth.
The width threshold $N_0$ is polynomial in
$d,n,L,\log(2/\beta)$, with degree depending only on $\zeta$.
\end{theorem}

The theorem concerns the Boolean function of a posterior draw. Its
low-product representative need not be the draw itself. In particular,
it would be incorrect to replace this result by an assertion that all
posterior parameters have small energy.

\begin{lemma}[Polynomial size of the ceiling]
\label{lem:ceiling-poly}
For $S=d+n+2$, the preceding product satisfies
\[
 P_0\le e^{19/2}544^{3/4}S^9.
\]
In particular its polynomial degree is independent of $L,\zeta,\beta$.
\end{lemma}
\begin{proof}
The displayed definitions give $C_0\le S$, $R_0\le4S^2$,
$M\le17S^7$, and $T\le272S^{12}$. Substitute these into $P_0$ and
use $d\ge1$. For example,
$\log(16(n+1))\le4S$ for $S\ge4$, which suffices for the bound on
$R_0$; the remaining inequalities follow directly.
\end{proof}

\section{From samples to an exact parallel circuit}
\label{sec:converse}

Let $\mathcal H_P$ be the finite set of Boolean functions on $X_d$
that admit globally unit-margin tanh representations of product at
most $P$, allowing arbitrary finite depths and widths. It is finite
because the domain is finite. Theorem~\ref{thm:main} gives uniform
circuit bounds for every member of this set.

\begin{theorem}[Sample-to-parallel-circuit converse]
\label{thm:converse}
Fix $d,n\ge1$, $L\ge3$, $\zeta>0$, and $0\le\varepsilon<1/2$.
Use Theorem~\ref{thm:ceiling} with $\beta=1/6$, and fix any one
$N\ge N_0$ and any one $0<\sigma^2\le N^{-(1+\zeta)}$.
These choices are independent of the input law.

Suppose that for every input distribution $\mathcal D$ on $X_d$,
an iid $n$-sample labeled by $f$ has positive probability of satisfying
\[
 \mathcal D\{x:p_S(f(x)\mid x)\le2/3\}\le\varepsilon.
\]
Put $\kappa=1-2\varepsilon>0$. Then $f$ has a threshold circuit of
size polynomial in $d+n+\kappa^{-1}$ and depth $O(\log(d+n+2))$.
It has a bounded-fan-in Boolean circuit of polynomial size in the same
quantities and depth
\[
 O\bigl(\log^2(d+n+2)+\log(\kappa^{-1}+2)\bigr).
\]
All polynomial degrees and depth constants are universal. They do
not depend on the original teacher, $L$, $N$, or $\zeta$.
\end{theorem}
\begin{proof}
Fix $\mathcal D$ and select a successful sample. The ceiling supplies
a posterior event of mass at least $5/6$ on which the represented
Boolean function belongs to $\mathcal H_{P_0}$. Condition on that
event, obtaining a distribution $\nu$ on this finite class. On the
good set
\[
 G=\{x:p_S(f(x)\mid x)>2/3\},
\]
whose $\mathcal D$-mass is at least $1-\varepsilon$, a draw from
$\nu$ has wrong-label probability strictly less than
$(1/3)/(5/6)=2/5$. Its signed vote therefore has expectation greater
than $1/5$ at every $x\in G$.

Choose an odd integer $m_0\ge100(d+1)$ no larger than $100(d+1)+2$.
Independent draws of $m_0$ functions from $\nu$ have majority error
at a fixed good input at most $e^{-m_0/50}$ by Hoeffding's inequality.
A union bound over $X_d$ is strictly less than one. Thus some majority
$g$ of $m_0$ members of $\mathcal H_{P_0}$ agrees with $f$ on all of
$G$, and has $\mathcal D$-error at most $\varepsilon$.

Let $\mathcal K$ be the finite class of all such $m_0$-voter
majorities. This class is independent of $\mathcal D$. The conclusion
just proved says
\[
 \text{for every }\mathcal D,\quad
 \max_{g\in\mathcal K}\E_{x\sim\mathcal D}[f(x)g(x)]\ge\kappa.
\]
The finite minimax theorem, applied to the matrix with entries
$f(x)g(x)$, gives a distribution $\lambda$ on $\mathcal K$ with
\[
 f(x)\E_{g\sim\lambda}g(x)\ge\kappa
 \qquad\text{for every }x\in X_d.
\]
Choose an odd integer $m_1\ge4(d+1)/\kappa^2$ within two of this
bound. Hoeffding and a union bound again give
\[
 2^d e^{-m_1\kappa^2/2}<1.
\]
Some majority of $m_1$ members of $\mathcal K$ therefore equals $f$
on every input. This is a majority of majorities of low-product
functions; it is not being asserted to have low product itself.

Compile all its constituent functions by Theorem~\ref{thm:main}.
Lemma~\ref{lem:ceiling-poly} bounds $P_0$ by a universal polynomial
in $d+n$. Put these circuits in parallel and add the two majority
layers. Their number is polynomial in $d+n+\kappa^{-1}$, and the
threshold depth increases only by two. For the Boolean circuit,
convert the constituent circuits first, then implement the two
majorities in depths $O(\log m_0)$ and $O(\log m_1)$ using the
construction in Section~\ref{sec:compiler}. This gives the sharper
additive $O(\log\kappa^{-1})$ term in the stated bound.
\end{proof}

\begin{remark}[Quantifiers and mixtures]
The hypothesis is distribution-free, and only existence of a
successful sample is used for each distribution. The samples, low-product
representatives, and both lists of voters may depend nonuniformly on
$f$. The one width and one actual Gaussian scale in the theorem are
fixed before the distribution is chosen. No posterior algorithm,
posterior concentration on a single function, or pointwise statement
about an individual posterior draw is assumed.
\end{remark}

\section{The critical endpoint: canonical features and parallel iteration}
\label{sec:critical-parallel}

The critical prior has $t=1$, equivalently $\sigma^2=1/N$. Its
finite-width rate ceiling is different from the product ceiling above.
We prove a separate circuit theorem for its near-baseline rate class.
The proof preserves a finite set of anchor correlations, rather than
the full feature covariance, and retains the learner depth in the
resulting resource bound.

\subsection{Variational paths and the imported scalar facts}

For a finite input set $X\subseteq X_d$, put
$Q_0(x,y)=(1+x\cdot y)/d$. A critical controlled path with
$H=L-1$ hidden transitions consists of laws
\[
 \rho_j\ll\mathcal N(0,Q_{j-1}),\qquad
 Q_j(x,y)=\E_{\rho_j}[\tanh(G_x)\tanh(G_y)],
 \qquad D_j=\KL(\rho_j\Vert\mathcal N(0,Q_{j-1})).
\]
Singular Gaussian references are allowed. A readout vector
$z\in\operatorname{range}(Q_H)$ has cost
\[
 r=\tfrac12z^TQ_H^\dagger z,
 \qquad A=\sum_{j=1}^H D_j+r.
\]
Equivalently, if $Q_H$ is the Gram matrix of feature vectors $u_H(x)$,
then $z(x)=\langle c,u_H(x)\rangle$ for a vector $c$ of squared norm
$2r$. Let $J_L(f)$ be the infimum of $A$ over full-cube paths with
$f(x)z(x)\ge1/2$ everywhere.

The scalar value $V(q)$ is the infimum of
\[
 \sum_{j=1}^m\KL(\eta_j\Vert\mathcal N(0,q_{j-1}))
       +\frac1{8q_m},\qquad
 q_0=q,\quad q_j=\E_{\eta_j}\tanh^2 G,
\]
over all finite lengths, including length zero. We use the following
previously proved facts from \cite[``Scalar control and target-component
constructions'']{critical}. They are stated explicitly to identify the
only scalar-control input to our new parallelization argument.

\begin{proposition}[Imported scalar Bellman facts]
\label{prop:scalar-facts}
There is a limit $V_*:=\lim_{q\downarrow0}V(q)\in(1,2)$ and universal
$C_0,c>0$ such that
\[
 0\le V_*-V(q)\le C_0q\quad(0<q\le2),\qquad
 V(q)\le\frac1{8q}.
\]
For a scalar transition from reference variance $p$ to feature variance
$q$, with entropy $D$, one has
\[
 D\ge V(p)-V(q),\qquad
 D\ge q-p+cp^2\quad(0<p\le2).
\]
The constant $c$ may be decreased so that $p-cp^2$ is increasing on
$[0,5]$.
\end{proposition}

The source states the small-state bound first and extends it by
boundedness; increasing $C_0$ gives the interval above. Its strict
potential proof gives dissipation on $[0,1]$. The same proof extends
to $[0,2]$: for $q>0$,
$\log\E\exp\{\tanh^2(\sqrt q Z)\}<q$, its difference from $q$ is
$q^2+O(q^3)$ at zero, and compactness supplies a positive uniform
constant. The Bellman inequality is also immediate by appending a
nearly optimal scalar suffix.

For $D\ge0$, define $\Gain(D)\ge1$ by
\[
 \Gain(D)-1-\log\Gain(D)=2D.
\]
Gaussian entropy duality gives, for any tilt $\rho$ of a standard
Gaussian with entropy at most $D$,
\begin{equation}
 \E_\rho\langle g,u\rangle^2\le\Gain(D)\norm u_2^2,
 \qquad \sqrt{\Gain(D)}\le e^{\sqrt D},
 \qquad \Gain(D)\le4D+2.
 \label{eq:critical-gain}
\end{equation}
Indeed, for a scalar second moment $m>1$, maximizing the entropy
dual over quadratic functions gives
$D\ge(m-1-\log m)/2$. For the second inequality write
$v=\tfrac12\log\Gain(D)$ and use
$D=(e^{2v}-1-2v)/2\ge v^2$. For the last, use
$\log u\le u/2$ for $u\ge2$. These facts also show that every
linear coordinate lies in $L^2(\rho)$.

\begin{lemma}[Variance and entropy occupation]
\label{lem:critical-occupation}
Let a full-cube path have global margin $1/2$ and action
$A\le V_*+\epsilon$, where $\epsilon\ge0$. Put
$q_j=\max_x Q_j(x,x)$. Then
\[
 \sum_{j=0}^{H-1}q_j+\sum_{j=1}^H\sqrt{D_j}
 \le C\bigl[\log(H+2)+\sqrt{H\epsilon}\bigr]
\]
for a universal $C$.
\end{lemma}
\begin{proof}
All one-input variances are positive: zero variance would propagate
at finite entropy and make the final output zero. The initial
diagonal is $(d+1)/d\le2$, and subsequent diagonals are at most one.
Project a transition onto an input attaining its new maximum, and
use the monotonicity of $p-cp^2$ in Proposition~\ref{prop:scalar-facts}.
This gives
\begin{equation}
 D_j\ge q_j-q_{j-1}+cq_{j-1}^2.
 \label{eq:critical-max-dissipation}
\end{equation}
At an input attaining $q_i$, scalar projection of the suffix and
Cauchy--Schwarz for the readout show that its remaining cost is at
least $V(q_i)$. Therefore, with $S_i=\sum_{j\le i}D_j$,
\begin{equation}
 S_i\le\epsilon+C_0q_i,
 \qquad S_H\le V_*+\epsilon.
 \label{eq:critical-prefix-budget}
\end{equation}

We give the occupation argument because both its conclusions will
be used. On a consecutive interval of $\ell$ transitions with
pre-states at least $u>0$, entropy expenditure $E$, and endpoint
also at least $u$, dividing~\eqref{eq:critical-max-dissipation} by
the squared pre-state and telescoping yields
\begin{equation}
 \ell\le C(E/u^2+1/u).
 \label{eq:critical-reciprocal-occupation}
\end{equation}
The algebra is
$(q'-q)/q^2\ge1/q-1/q'$, with difference
$(q'-q)^2/(q^2q')\ge0$.

We also need a bound for an initial prefix of length $t$ and entropy
expenditure $E<3$:
\begin{equation}
 \sum_{j<t}q_j\le C\log(t+2)+C\sqrt{tE}.
 \label{eq:critical-forced-cooling}
\end{equation}
Let $F(v)=v-cv^2$. Compare $q_j$ with $r_j+s_j$, where
$r_0=q_0$, $s_0=0$, $r_{j+1}=F(r_j)$ and
$s_{j+1}=F(s_j)+D_{j+1}$. These sequences are nonnegative,
$r_j\le2$, $s_j\le E<3$, and $F$ is increasing on their range.
Since $F(r+s)\le F(r)+F(s)$, induction gives $q_j\le r_j+s_j$.
Also $r_j\le1/(cj+1/r_0)$ and
$c\sum_{j<t}s_j^2\le E$. Summing and applying Cauchy--Schwarz
proves~\eqref{eq:critical-forced-cooling}.

If $\epsilon\ge1/4$, summed dissipation and Cauchy--Schwarz give
both claims directly from $S_H\le2+\epsilon$. Otherwise $S_H<3$.
Put $a=\max\{\epsilon,(H+1)^{-1}\}$ and $b=\min_j q_j$.
If $a\ge1/4$, then $H$ is bounded; if $b\ge1/4$,
\eqref{eq:critical-reciprocal-occupation} bounds $H$. Thus assume
$u_0=\max\{a,b\}<1/4$. Let $u_j=2^ju_0$, stop at the first
$u_J\ge1/4$, and define last-exit indices
$\tau_j=\max\{i:q_i\le u_j\}$. There are
$J=O(\log(H+2))$ bands.

For the initial prefix through $\tau_0$, if $b\le a$ its entropy
is $O(\epsilon+1/H)$ and its length at most $H$.
Equation~\eqref{eq:critical-forced-cooling} gives its state sum,
and Cauchy--Schwarz gives its square-root entropy sum, bounded by
$C[\log(H+2)+\sqrt{H\epsilon}]$. If $b>a$, its endpoint is $b$,
its entropy is $O(b)$, and all its states are at least $b$.
Equation~\eqref{eq:critical-reciprocal-occupation} gives length
$O(1/b)$; forced cooling bounds its state sum by $O(\log(H+2))$,
and Cauchy--Schwarz bounds its square-root entropy by a constant.

For a nonempty later block $(\tau_j,\tau_{j+1}]$, its entropy is
at most $\epsilon+2C_0u_j=O(u_j)$. After deleting its first
transition, all pre-states and its endpoint exceed $u_j$, so its
length is $O(1/u_j)$. Restoring the first transition preserves this
bound. Its initial state is at most $u_j$, and each subsequent
increase is bounded by the entropy spent; all states in the block
are therefore $O(u_j)$. The block's state sum and square-root
entropy sum are both $O(1)$. The tail after $\tau_J$ has bounded
length and bounded contributions by the same occupation bound.
Summing the blocks proves the lemma.
\end{proof}

\subsection{Anchors and a finite canonical state}

For two positive semidefinite kernels with identical anchor block
$K_{TT}=K'_{TT}$, define
\[
 d_T(K,K')=\max_{x\in X_d}
 \sqrt{\norm{a_x-a'_x}_2^2+(\sqrt{v_x}-\sqrt{v'_x})^2},
\]
where $a_x=K_{TT}^{\dagger/2}K_{Tx}$ and
$v_x=K_{xx}-\norm{a_x}_2^2$, and similarly for $K'$.
This has a geometric meaning: identify the anchor spans isometrically,
project each query onto that common span, and align its remaining
orthogonal direction. It is also the root mean square difference in
a Gaussian coupling that shares the anchor vector and the one-query
residual Gaussian. Any coupling of query feature vectors with identical
anchor vectors bounds this distance by its query root mean square
error. No inverse Gram matrix will be evaluated by the final circuit.

\begin{lemma}[Finite anchor replacement]
\label{lem:critical-anchors}
Let a full-cube controlled path have action $A\le V_*+\epsilon$,
and put $T_0=1+\log(H+2)+\sqrt{H\epsilon}$.
There is a genuine controlled path whose row densities depend only
on a common set of
\[
 q\le\poly(d,H,A+1)\exp(CT_0)
\]
anchors, with the same anchor covariance blocks as the original
path. It has an anchor-supported readout of global margin $1/2$
and action at most $A+C/H$. Its layer entropies are at most the
original $D_j$. Its readout vector has norm at most $C\sqrt{A+1}$.
\end{lemma}
\begin{proof}
Write $I_j(T)$ for the entropy of the layer-$j$ anchor marginal
relative to its Gaussian reference. It is monotone in $T$ and at
most $D_j$. Choose $0<\rho\le1/64$ of order
$[H(A+1)]^{-1}$, and a cutoff $\delta>0$ specified below.
Greedily add an input whenever either
$I_j(T\cup\{x\})-I_j(T)>\delta$ for some $j$, or projection of
the final readout onto the anchor feature span changes its output
at $x$ by more than $\rho$. There are at most
$A/\delta$ entropy additions. A readout addition increases the
squared norm of its projection by at least $\rho^2$, since the
final feature vectors have norm at most one. There are at most
$2r/\rho^2$ such additions. This proves termination and a bound
on the common anchor count.

At layer $j$ replace the original density by its conditional
expectation given $G_T$, denoted $w_j(G_T)$. This preserves the
anchor marginal exactly and has entropy $I_j(T)\le D_j$.
For a query $x$, the conditional entropy of $G_x$ beyond $G_T$
is at most $\delta$. Conditional maximal coupling and Pinsker's
inequality give a coupling of the original and replaced bounded
tanh features, with identical anchor features, whose squared query
error is at most $4\sqrt{\delta/2}$. Indeed the mismatch probability
is conditional total variation, whose average is at most
$\sqrt{\delta/2}$, and the squared feature difference is at most
four. Thus the one-step $d_T$ error is at most $C\delta^{1/4}$.
Pinsker's inequality here follows from binary entropy after data
processing onto a set attaining total variation.

Transfer the anchor density $w_j$ to any preceding kernel with the
same anchor block. Couple its Gaussian anchor-plus-query vectors
as in the definition of $d_T$. The common density has entropy at
most $D_j$ relative to this joint Gaussian law. Equation
\eqref{eq:critical-gain} and tanh's Lipschitz property imply that
the resulting $d_T$ error is at most
$\sqrt{\Gain(D_j)}$ times the input error. Applying these
replacements successively defines a genuine full-cube controlled
path, with exact original anchor blocks. Its total query error is
at most
\[
 C H\delta^{1/4}\exp\left(\sum_j\sqrt{D_j}\right).
\]
Choose $\delta$ so this is at most
$\rho/\sqrt{\max\{2r,1\}}$.
Lemma~\ref{lem:critical-occupation} makes $\delta^{-1}$ at most a
polynomial in $H,A+1$ times $\exp(CT_0)$, as needed for the anchor
count.

The projected readout has norm at most $\sqrt{2r}$. Its query
values change by at most $\rho$ during the path transfer, in
addition to the projection error $\rho$. Thus it has margin at
least $1/2-2\rho$. Scale it by $(1-4\rho)^{-1}$ to restore
margin $1/2$. This changes its action by $O(\rho A)=O(1/H)$ and
preserves the asserted norm bound. The row entropy costs do not
increase.
\end{proof}

\begin{lemma}[Canonicalization with additive dimension]
\label{lem:critical-canonical}
Let the query states $u_x\in\R^K$ include anchor states $u_a$.
Let $\rho$ be a tilt of $g\sim\mathcal N(0,I_K)$ of entropy at
most $D$, and set $\Phi(u)=\tanh(\langle g,u\rangle)$ as a vector
in $\mathcal H=L^2(\rho)$. There is a map
$F:\R^K\to\R^{K+q+1}$, allowing unused zero coordinates, such that
its Gram kernel has the same anchor block, all query-anchor inner
products, and all query norms as $\Phi$. Moreover
\[
 F(u)=Au+B(u),\qquad
 \norm A_{\mathrm{op}}\le\sqrt{\Gain(D)},\qquad B(0)=0,
\]
and on the ball $\norm u\le\sqrt Q$,
\[
 \Lip(B)\le C(Q+\sqrt D).
\]
No assertion of full query-query kernel preservation is made.
\end{lemma}
\begin{proof}
Let $V\subseteq\mathcal H$ be the span of the $K$ functions $g_i$
and the $q$ anchor functions $\Phi(u_a)$. For an orthonormal basis
$e_1,\ldots,e_M$ of $V$, with $M\le K+q$, define
\[
 F(u)=\bigl(\langle e_1,\Phi(u)\rangle,\ldots,
       \langle e_M,\Phi(u)\rangle,
       \norm{(I-P_V)\Phi(u)}\bigr).
\]
Pythagoras preserves every norm. At an anchor the last coordinate
is zero, which preserves every query-anchor inner product and the
anchor block. Since $L(u)=\langle g,u\rangle\in V$, set
\[
 Au=(\langle e_b,L(u)\rangle)_{b=1}^M\oplus0.
\]
Equation~\eqref{eq:critical-gain} bounds its operator norm. Writing
$R(u)=\Phi(u)-L(u)$, the residual is
$B(u)=(P_VR(u),\norm{(I-P_V)R(u)})$ in these coordinates.
Orthogonality and the reverse triangle inequality give
\[
 \norm{B(u)-B(v)}^2\le\norm{R(u)-R(v)}_{L^2(\rho)}^2.
\]

For the last bound put $X=\langle g,u\rangle$,
$Y=\langle g,v\rangle$ and
$Z=\langle g,u-v\rangle/\norm{u-v}$ when $u\ne v$.
The derivative of $\tanh s-s$ has absolute value $\tanh^2s$, so
\[
 \frac{|R(u)-R(v)|^2}{\norm{u-v}^2}
 \le W:=Z^2\min\{1,\max(|X|,|Y|)^4\}.
\]
Under the reference Gaussian, $Z$ is standard normal and
$\operatorname{Var}(X),\operatorname{Var}(Y)\le Q$. For a fixed
sufficiently small $c>0$,
\[
 e^{cW}-1\le cZ^2(X^4+Y^4)e^{cZ^2}.
\]
Conditioning $X$ or $Y$ on $Z$ bounds its fourth moment by
$CQ^2(1+Z^4)$, regardless of their correlations. Hence
$\log\E e^{cW}\le CQ^2$. Entropy duality gives
$\E_\rho W\le C(D+Q^2)$, proving the Lipschitz assertion.
\end{proof}

Apply this canonicalization successively to the path from
Lemma~\ref{lem:critical-anchors}. At each stage the reference
Gaussian can be written in the current finite coordinates, tilted
by the prescribed anchor density. The next anchor-plus-one-query
Gaussian law depends only on the three types of moments preserved
by Lemma~\ref{lem:critical-canonical}. Induction therefore gives
$d_T=0$ between the canonical recursion and the genuine anchor-only
path. In particular their diagonals and anchored outputs are equal.
After $H$ transitions the canonical dimension is at most
$d+1+H(q+1)$. This does not claim that the canonical full covariance
itself is the original controlled transition covariance.

\subsection{Uniform finite evaluation of a canonical map}

\begin{lemma}[Polynomial canonical evaluators]
\label{lem:critical-evaluation}
In Lemma~\ref{lem:critical-canonical}, suppose
$K,q,\Gain(D),U,\eta^{-1}$ are bounded by a fixed polynomial in
$R\ge2$. On the ball $\norm u\le U$, the map $F$, and the residual
$F(u)-Au$, admit uniform approximations of error $\eta$ by
polynomial-size constant-depth threshold circuits on a sufficiently
fine inverse-polynomial input grid. The circuits use bounded dyadic
constants of $O(\log R)$ bits. The polynomial bounds depend only
on the stipulated resource polynomial, not on an anchor condition
number or a description of $\rho$.
\end{lemma}
\begin{proof}
Write
\[
 a(u)=\E_\rho[e(g)\tanh(g\cdot u)],\quad
 b(u)=\E_\rho\tanh^2(g\cdot u),\quad
 F(u)=(a(u),\sqrt{b(u)-\norm{a(u)}^2}).
\]
The radicand is nonnegative and $\norm{a(u)}\le1$.
By~\eqref{eq:critical-gain},
$\E_\rho\norm g^2\le K\Gain(D)$. Truncate to $\norm g\le B$.
The discarded probability is at most $K\Gain(D)/B^2$. Bessel's
inequality bounds the discarded vector expectation uniformly in $u$
by the square root of that probability; the discarded scalar
expectation is bounded by the probability itself. A polynomial $B$
makes both errors any prescribed inverse-polynomial amount.

On the retained set use the importance density
\[
 p(g)=\frac{(1+\norm{e(g)}^2)\boldsymbol1_{\{\norm g\le B\}}}{Z},
 \qquad Z\le M+1,
\]
relative to $\rho$. The retained set has positive probability.
Its vector and scalar integrands have respective bounds
$\norm{e(g)\tanh(g\cdot u)/p(g)}\le Z/2$ and
$\tanh^2(g\cdot u)/p(g)\le Z$.
They have Lipschitz bounds at most $ZB/2$ and $2ZB$.
A net of the radius-$U$ input ball at spacing proportional to
$\theta/(ZB)$ has logarithmic cardinality
$O(K\log(2+UZB/\theta))$. The Hilbert-valued bounded-differences
argument used in Lemma~\ref{lem:maurey}, and a union bound on this
net, give a polynomial number of importance samples uniformly
approximating $a,b$ to error $\theta$. Lipschitzness extends the
bound to the whole ball.

The formulas have the form
\[
 \widehat a(u)=\sum_{\ell=1}^m c_\ell\tanh(g_\ell\cdot u),
 \qquad
 \widehat b(u)=\sum_{\ell=1}^m b_\ell\tanh^2(g_\ell\cdot u),
\]
where $\norm{g_\ell}\le B$, $\norm{c_\ell}\le Z/(2m)$ and
$0\le b_\ell\le Z/m$. In particular, large coefficients in a
possibly ill-conditioned formula for the basis $e$ never appear in
the evaluator. If both errors are at most $\theta\le1$, then
\[
 \left|\widehat b-\norm{\widehat a}^2
                  -(b-\norm a^2)\right|\le3\theta+\theta^2.
\]
Clip the approximate radicand at zero. The inequality
$|\sqrt s-\sqrt t|\le\sqrt{|s-t|}$ shows that
$\theta$ of order $\eta^2$ suffices for error $\eta$ in $F$.

All incoming nodes and output coefficients are polynomially bounded
and can be rounded on inverse-polynomial meshes. The matrix $A$ is
an independent nonuniform constant of bounded operator norm and
can be rounded directly; it need not be estimated from truncated
samples. Affine forms, tanh lookups, squares, finite sums, clipping,
and square roots now have polynomial range and inverse-polynomial
precision requirements. Use the constant-threshold-depth lookup
and comparison construction from Section~\ref{sec:compiler} for
each. Square root needs the squared accuracy mesh. This proves the
claim. It uses existence of finite quadrature constants, not an
algorithm for the density or its orthonormal basis.
\end{proof}

\begin{theorem}[Near-baseline critical rate implies parallel circuits]
\label{thm:critical-normal-form}
Suppose $f$ has a globally $1/2$-margined critical path with $H$ hidden
transitions and action $A\le V_*+\epsilon$, $\epsilon\ge0$.
Put
\[
 T_0=1+\log(H+2)+\sqrt{H\epsilon},\qquad
 R=(d+H+A+2)e^{CT_0},
\]
with a sufficiently large universal $C$.
Then $f$ has a threshold circuit of size $R^{O(1)}$ and depth
$O(\log R)$, and a bounded-fan-in Boolean circuit of size $R^{O(1)}$
and depth $O(\log^2R)$, with universal constants.
Consequently, if $d+L\le Q^k$ and
$J_L(f)\le V_*+k\log^2Q/L$ for fixed $k$, then $f$ has size
polynomial in $Q$ and threshold depth $O_k(\log Q)$.
\end{theorem}
\begin{proof}
Use Lemma~\ref{lem:critical-anchors} and the canonical recursion.
Its dimension is polynomial in $R$. The genuine anchor-only path
has action at most $A+C/H$, the same diagonals and readout as the
canonical recursion, and row entropies at most $D_j$.
Apply Lemma~\ref{lem:critical-occupation} to that genuine path for
its diagonal sum, with the additional $C/H$ absorbed in the constant,
and to the original path for its original entropy sum. Together they give
\[
 \sum_j Q_{j-1}^{\max}+\sum_j\sqrt{D_j}\le CT_0.
\]
Write its canonical maps as $F_j(u)=A_ju+B_j(u)$, with
\[
 \alpha_j=\sqrt{\Gain(D_j)},\quad
 \norm{A_j}_{\mathrm{op}}\le\alpha_j,\quad
 \Lip(B_j)\le\ell_j:=C(Q_{j-1}^{\max}+\sqrt{D_j})
\]
on the relevant state balls. Thus
$\sum_j\ell_j\le CT_0$ and $\prod_j\alpha_j\le e^{CT_0}$.

Extend each $B_j$ by radial projection onto its state ball. For
finite evaluation, first floor each radius at $R^{-c}$ and then round
it upward on a dyadic mesh of spacing at most $R^{-c}$, for a
sufficiently large universal $c$. The increase in the sum of squared
radii is $O(HR^{-c})\le O(1)$, so the total Lipschitz budget increases
by at most a constant. The actual path is unchanged.
These extensions fix zero, remain Lipschitz, and have polynomially
bounded ranges. Radial projection is uniformly evaluable with
polynomial resources by norm, square-root, and reciprocal lookups;
division only occurs above the inverse-polynomial radius floor.

Let $S_j=\prod_{i\le j}\alpha_i$ and $v_j=u_j/S_j$, with $S_0=1$.
The normalized recursion is
\[
 v_j=C_jv_{j-1}+H_j(v_{j-1}),\qquad
 C_j=A_j/\alpha_j,\quad \norm{C_j}_{\mathrm{op}}\le1,
\]
where $H_j(v)=B_j(S_{j-1}v)/S_j$ uses the extended residual.
Its Lipschitz constant is at most $\ell_j$. All $S_j$ lie between
one and a polynomial in $R$.

Precompute the contraction products
$C_{j:i}=C_j\cdots C_i$ as nonuniform matrices. Their entries have
absolute value at most one and their dimensions are polynomial.
The recursion is equivalent to the Volterra system
\[
 v_j=C_{j:1}v_0+\sum_{i=1}^j C_{j:i+1}H_i(v_{i-1}),
\]
where an empty product is the identity. Initialize
$v_j^{(0)}=C_{j:1}v_0$ and compute all times in parallel by
\begin{equation}
 v_j^{(k+1)}=C_{j:1}v_0+
             \sum_{i=1}^j C_{j:i+1}H_i(v_{i-1}^{(k)}).
 \label{eq:critical-picard}
\end{equation}
Put $\Lambda=\sum_i\ell_i=O(T_0)$. Expanding the error recurrence
over strictly increasing time indices gives
\[
 \max_j\norm{v_j^{(k)}-v_j}
 \le2\sqrt2\,\frac{\Lambda^k}{k!}.
\]
Indeed each product uses distinct ordered indices, and their sum
is at most $\Lambda^k/k!$. Thus
$k=O(\Lambda+\log(1/\eta))=O(\log R)$ iterations suffice for
any fixed inverse-polynomial error $\eta$.

Uniform additive error $\delta$ per update contributes at most
\[
 \delta\sum_{m\ge0}\frac{\Lambda^m}{m!}
 =\delta e^\Lambda=\delta\,R^{O(1)}.
\]
This treats finite quadrature and rounded arithmetic as errors
relative to the exact Lipschitz maps; the approximating maps need
not themselves have small Lipschitz constants. Clipped residuals
are bounded, so the right side of~\eqref{eq:critical-picard} has
polynomial norm at every iteration. This supplies a fixed polynomial
evaluation range and a bootstrap for its inexact versions. Divide
the primitive error tolerance by a sufficiently large polynomial
to absorb all matrix, coefficient, and summation errors in $\delta$.

Lemma~\ref{lem:critical-evaluation} makes every simultaneous
update a polynomial-size constant-depth threshold circuit. There
are polynomially many time indices and matrix entries, and
$O(\log R)$ iterations. The final rescaling by $S_H$ and bounded
readout amplify error only polynomially. Choose $\eta,\delta$ so
that the final output changes by less than $1/16$; its $1/2$ margin
then gives the exact signs. This proves the threshold bound.
The majority-gate simulation in Section~\ref{sec:compiler} gives
the Boolean bound. For an infimum $J_L(f)$, choose a path within
$1/H$ of that infimum. This adds only a constant to $T_0$, proving
the final assertion without an attainment assumption.
\end{proof}

\subsection{A finite-width full-cube rate ceiling}

We next prove the extraction step needed to apply the structural
theorem to the actual posterior, rather than a limiting rate model.

\begin{lemma}[Finite-event covariance certificate]
\label{lem:critical-dual-certificate}
Consider the critical width-$N$ prior on a set of $m$ inputs, with
$H$ hidden layers, and put $k=m(m+1)/2$. For any measurable raw
event $E$ of probability at least $e^{-NA}$, $A\ge0$, any
$\tau>0$, and $0<\Delta\le1$, there is a realization in $E$ and
laws $\rho_j\ll\mathcal N(0,Q_{j-1})$, relative to its empirical
preceding covariances, such that
\[
 \max_{x,y}\left|\E_{\rho_j}\tanh G_x\tanh G_y-Q_j(x,y)\right|
 \le\Delta,\qquad
 \sum_j\KL(\rho_j\Vert\mathcal N(0,Q_{j-1}))+r\le A+\tau,
\]
provided
\[
 N\ge\frac C\tau\left[
 Hk\log\left(1+\frac{C(A/\tau+1)Hk}{\Delta}\right)
 +m\log\bigl(C(A/\tau+1)\bigr)+1\right].
\]
Here $z$ is the selected realization's actual output and
$r=\tfrac12z^TQ_H^\dagger z$. No exchangeability of $E$ is required.
\end{lemma}
\begin{proof}
For a row law $\mu$ on $h\in[-1,1]^m$, let $\phi(h)$ list its $k$
upper-triangular products. Define the bounded dual
\[
 I_B(q;\mu)=\max_{\lambda\in[-B,B]^k}
 \{\lambda\cdot q-\log\E_\mu e^{\lambda\cdot\phi}\}.
\]
A maximizer exists. Its exponential tilt $\nu$ has moments $a$,
and coordinatewise optimality gives
$\lambda_i(q_i-a_i)=B|q_i-a_i|$. Consequently
\begin{equation}
 I_B(q;\mu)=\KL(\nu\Vert\mu)+B\norm{q-a}_1\ge0.
 \label{eq:critical-bounded-dual}
\end{equation}
The same feature tilt of the Gaussian preactivation has identical
entropy and moments.

For $N$ iid rows let $\widehat q$ be their empirical moments. A
grid $\Gamma\subset[-B,B]^k$ of covering radius $\eta$ has size
at most $(1+CB/\eta)^k$. The dual objective is $2k$-Lipschitz in
$\lambda$, and each fixed exponential likelihood ratio has mean
one. Hence
\[
 \E e^{N I_B(\widehat q;\mu)}
 \le e^{2Nk\eta}|\Gamma|.
\]
This bound is uniform in the conditional reference row law.
Iterate it forward through the network. Given the hidden layers,
the readout has covariance $Q_H/N$, so, for $0<s<1$,
$\E[e^{(1-s)Nr}\mid\text{hidden}]
=s^{-\operatorname{rank}(Q_H)/2}\le s^{-m/2}$.
Writing $I_j$ for the realized layer dual costs gives
\[
 \E_{\mathrm{prior}}e^{N[\sum_j I_j+(1-s)r]}
 \le e^{2NHk\eta}|\Gamma|^H s^{-m/2}.
\]
Conditioning on $E$ has relative entropy at most $NA$. Entropy
duality therefore bounds the conditional expectation of
$\sum I_j+(1-s)r$ by
\[
 A+2Hk\eta+
 \frac{H\log|\Gamma|+(m/2)\log(1/s)}N.
\]
Take $s=\tau/[16(A+\tau)]$,
$\eta=\tau/(64Hk)$, and $B=(A+\tau)/\Delta$.
Under the stated width bound some realization in $E$ has the last
cost at most $A+\tau/8$. Since all $I_j$ are nonnegative and
$(A+\tau)(1-s)=A+15\tau/16$, its unweighted total
$\sum I_j+r$ is at most $A+\tau$. Apply
\eqref{eq:critical-bounded-dual} layer by layer. Each moment
$\ell_1$ error is at most $I_j/B\le\Delta$, proving the lemma.
\end{proof}

\begin{lemma}[Consistent-path transfer]
\label{lem:critical-consistent-transfer}
For a certificate as above with total cost at most $B$, there is
a genuine controlled path of cost at most $B$, with output
$\widehat z$, satisfying
\[
 \max_x|\widehat z(x)-z(x)|
 \le \sqrt{2B}\,H\sqrt{m\Delta}\,e^{\sqrt{HB}}.
\]
\end{lemma}
\begin{proof}
Two positive semidefinite $m$-point kernels differing entrywise
by at most $\Delta$ have feature embeddings with each corresponding
row distance at most $\sqrt{m\Delta}$. To see this, their operator
norm difference is at most $m\Delta$, and
$\norm{K^{1/2}-{K'}^{1/2}}_{\mathrm{op}}
\le\sqrt{\norm{K-K'}_{\mathrm{op}}}$.
The latter follows by monotonicity of the matrix square root from
$K\le K'+aI$ and $K'\le K+aI$. Any existing embedding can be
identified isometrically with its square-root embedding, enlarging
the ambient Hilbert space when necessary.

Inductively couple Gaussian rows for the genuine preceding kernel
and the empirical reference using such common feature embeddings.
Tilt the joint Gaussian by the certificate's original density.
The new marginal has no larger entropy by data processing, while
\eqref{eq:critical-gain} bounds each root mean square preactivation
difference by $\sqrt{\Gain(D_j)}$ times its previous feature
difference. Apply tanh, then use the certificate moment error to
identify the controlled moment embedding with the next empirical
embedding at additional error $\sqrt{m\Delta}$. Couplings can be
joined by identifying the common feature span and adjoining their
orthogonal complements. Thus the maximal row error satisfies
\[
 e_j\le\sqrt{\Gain(D_j)}e_{j-1}+\sqrt{m\Delta},
 \qquad e_0=0.
\]
It follows that
$e_H\le H\sqrt{m\Delta}\exp(\sum_j\sqrt{D_j})
\le H\sqrt{m\Delta}e^{\sqrt{HB}}$.
Transfer the empirical readout vector of norm $\sqrt{2r}$ through
this common embedding and project it onto the new final feature
span. Its norm and readout cost do not increase, while its output
error is at most $\sqrt{2r}\,e_H\le\sqrt{2B}\,e_H$.
\end{proof}

For the raw mass of a sample memorizer we use the established
finite-width Fourier-seed result
\cite[``Finite-width target mass for Fourier seeds'']{critical}:
there are universal constants such that a Boolean target $g$ of
Fourier variation $A_F(g)$ has raw critical probability at least
\[
 \exp\left[-N\left(V_*+C\frac{\log^2 R_g}{L}\right)\right],
 \qquad R_g=(d+L+2)A_F(g),
\]
at every $N\ge R_g^C$ whenever $L\ge C\log R_g$.
The event here has the original strict global margin and the signs
of $g$. This is an imported finite-width target-mass theorem, not
a rate-limit assertion.

\begin{theorem}[Critical full-cube posterior ceiling]
\label{thm:critical-ceiling}
There are universal $C,c>0$ such that the following holds. Put
$R_0=d+n+L+2$, with $d,n\ge1$, and suppose
$L\ge C\log R_0$. At the actual critical scale $\sigma^2=1/N$,
for every
\[
 N\ge 2^{Cd}R_0^C+CL\log(1/\beta),\qquad0<\beta<1,
\]
and every consistently labeled sample of at most $n$ inputs, the
original posterior assigns probability at least $1-\beta$ to the
fixed full-cube class
\[
 \mathcal G=\left\{g:J_L(g)\le
             V_*+c\frac{\log^2R_0}{L}\right\}.
\]
All its members have threshold circuits of polynomial size in
$R_0$ and depth $O(\log R_0)$, uniformly over the class.
\end{theorem}
\begin{proof}
If $x_1,\ldots,x_m$ are the distinct negative sample inputs, the
function $g_S=1-2\sum_{i=1}^m\boldsymbol1_{\{x_i\}}$ has the
observed signs and Fourier variation at most $1+2n$. Indeed each
point indicator has Fourier variation one. The imported target
mass theorem gives a lower bound $e^{-Na_0}$ for the conditioning
denominator, where
$a_0=V_*+C_0\log^2R_0/L$, at every $N\ge R_0^{C_0}$ after
enlarging constants.

Let $E$ be the raw event of globally margined draws whose Boolean
function $g$ has $J_L(g)>a_0+K/L$, where $K$ is a sufficiently
large universal constant. It is measurable because there are
finitely many sign functions. Suppose its probability were at
least $e^{-N(a_0+1/L)}$. Apply
Lemma~\ref{lem:critical-dual-certificate} on the entire cube, with
$m=2^d$, $\tau=1/L$, and action parameter $a_0+1/L$.
Choose $\Delta$ sufficiently small that
Lemma~\ref{lem:critical-consistent-transfer} gives output error
at most $\rho=[C L(a_0+1)]^{-1}<1/4$.
It suffices to choose a tolerance with
\[
 \log(1/\Delta)
 \le C\bigl[d+\log(L(a_0+2))+\sqrt{L(a_0+2)}\bigr].
\]
The certificate's required width is thus at most
$2^{Cd}R_0^C$. The transferred path preserves the selected
realization's signs; multiplying its readout by $(1-2\rho)^{-1}$
restores margin $1/2$ and adds at most $C/L$ action. It gives a
path of action at most $a_0+K'/L$ for one of the functions in $E$,
contradicting $K>K'$.

Hence the bad raw mass is less than $e^{-N(a_0+1/L)}$.
Divide by the sample-consistent denominator $e^{-Na_0}$ to bound
its posterior mass by $e^{-N/L}\le\beta$. Absorb $K/L$ into the
displayed logarithmic excess. Theorem~\ref{thm:critical-normal-form}
then gives the uniform circuit bounds. Its use of an arbitrarily
near-minimizing path removes any issue of rate attainment.
\end{proof}

\begin{corollary}[Critical sample-to-parallel-circuit converse]
\label{cor:critical-converse}
Let $0\le\varepsilon<1/2$, $\kappa=1-2\varepsilon$, and
$R_0=d+n+L+2$. Suppose $L\ge C\log R_0$ and choose one width
above the threshold of Theorem~\ref{thm:critical-ceiling} with
$\beta=1/6$. If, at $\sigma^2=1/N$, every input law gives
positive probability to an $n$-sample with confidence-failure
mass at most $\varepsilon$, then $f$ has threshold size polynomial
in $R_0+\kappa^{-1}$ and depth $O(\log R_0)$. Its Boolean depth
is $O(\log^2R_0+\log(\kappa^{-1}+2))$, with polynomial size.
\end{corollary}
\begin{proof}
Apply the two-majority argument of Theorem~\ref{thm:converse}
with the fixed class $\mathcal G$ from
Theorem~\ref{thm:critical-ceiling} in place of
$\mathcal H_{P_0}$. The one width is chosen before the input law,
and the ceiling is simultaneous over every realized sample.
Each voter now has the uniform circuit bound from
Theorem~\ref{thm:critical-normal-form}. No other step changes.
\end{proof}

The exponential width threshold in this converse is deliberate:
we extract on all $2^d$ inputs to obtain a distribution-independent
full-cube circuit class. It does not assert a polynomial width
threshold for this endpoint converse. It suffices for a necessary
condition on a sample guarantee required at every admissible width.

\section{Every fixed subcritical prefactor}
\label{sec:fixed-prefactor}

Fix $0<t<1$ while allowing the width to grow. The corresponding
exponent is $\gamma_N=1+2\log(1/t)/\log N$, so a fixed-gap theorem
does not apply uniformly in $N$. We prove a separate converse.
All constants and polynomial exponents in this section may depend
on the fixed value of $t$. Put
\[
 a=t^{-2}-1,\quad c_t=\frac{a-\log(1+a)}2,\quad
 \delta_t=\frac{(a+1)\log(1+a)-(a+2)\log(1+a/2)}2.
\]
A controlled path uses reference $\mathcal N(0,t^2Q_{j-1})$,
row entropies $D_j$, and readout action
\[
 r_t=\frac1{2t^2}z^TQ_H^\dagger z,\qquad
 A=\sum_{j=1}^H D_j+r_t,\qquad H=L-1.
\]
Its full-cube rate $J_{t,L}(f)$ is the infimum over such paths
with signed output margin at least $1/2$.

\subsection{A strict spectral covering gain}

\begin{lemma}[Entropy and two principal directions]
\label{lem:prefactor-spectral}
Let $\rho\ll\mathcal N(0,t^2I_K)$ have entropy at most $D$,
where $K\ge2$, and let $\lambda_1\ge\lambda_2$ be the largest
eigenvalues of $\Sigma=\E_\rho gg^T$. Set
$\alpha=\sqrt{\lambda_1}$ and
$b=(\lambda_1+\lambda_2)/(2\sqrt{\lambda_1})$. Then
\[
 \log\alpha\le(D-c_t)/a,\qquad
 \log b\le(D-c_t-\delta_t)/a,\qquad \delta_t>0.
\]
Every map satisfying
$\norm{F(u)-F(v)}\le\norm{\Sigma^{1/2}(u-v)}$
sends any ball of radius $r$ into two balls of radius $rb$.
\end{lemma}
\begin{proof}
Finite Gaussian entropy implies finite second moments, and absolute
continuity makes $\Sigma$ positive definite. Define
$F_t(\lambda)=(\lambda/t^2-1-\log(\lambda/t^2))/2$.
The exact identity
\[
 \KL(\rho\Vert\mathcal N(0,t^2I))
 =\KL(\rho\Vert\mathcal N(0,\Sigma))+\sum_iF_t(\lambda_i)
\]
uses the full second moment and allows a nonzero mean.
Thus $F_t(\lambda_1)+F_t(\lambda_2)\le D$.
The first gain bound follows from
\[
 F_t(\lambda)-c_t-\frac a2\log\lambda
 =\frac{a+1}2(\lambda-1-\log\lambda)\ge0.
\]
For the second put $s=\lambda_2/\lambda_1\le1$,
$u=(a+1)(\lambda_1+\lambda_2)/(a+2)$ and $v=(a+2)/2$.
Direct expansion gives
\begin{align*}
 F_t(\lambda_1)+F_t(\lambda_2)-c_t-a\log b-\delta_t
 ={}&v(u-1-\log u)\\
 &+\log\frac{1+s}{2\sqrt s}+\frac a2\log\frac2{1+s}\ge0.
\end{align*}
The signs follow from $\log u\le u-1$, arithmetic--geometric
mean, and $s\le1$. Also $\delta_t=0$ at $a=0$, and its derivative
is $\tfrac12\log((1+a)/(1+a/2))>0$ for $a>0$.

For the covering statement, take two centers on the top principal
axis of $\Sigma^{1/2}B_1$, at displacements
$\pm(\lambda_1-\lambda_2)/(2\sqrt{\lambda_1})$.
At top input coordinate of magnitude $s\in[0,1]$, the squared
distance to the nearer center is at most
\[
 \lambda_2+(\lambda_1-\lambda_2)(s^2-s)
                   +\frac{(\lambda_1-\lambda_2)^2}{4\lambda_1}
 \le b^2.
\]
The centers have preimages inside the input ball. Map these
preimages through $F$, then translate and scale the ball.
\end{proof}

The canonical map from Lemma~\ref{lem:critical-canonical}
has this metric property at every prefactor: the feature map
$\Phi(u)=\tanh(\langle g,u\rangle)$ does in $L^2(\rho)$,
and projection together with the residual norm is a contraction.
All its anchor and diagonal identities remain valid.
The gap $\delta_t\sim a^2/8$ as $t\to1$ is positive only away
from the endpoint.

\subsection{Normalized anchors and finite evaluation}

\begin{lemma}[Normalized anchor replacement]
\label{lem:prefactor-anchors}
Suppose a full-cube path has signed margin $1/2$ and
$A\le Hc_t+E$, $E\ge0$. Put $R=(d+H+A+2)e^E$.
There is a genuine path whose densities depend on a common set
of $R^{O_t(1)}$ anchors, with exact original anchor blocks and
an anchor-supported readout of margin $1/2$ and action at most
$A+C/H$. Its row entropies are at most the original $D_j$.
Its canonical state has dimension $R^{O_t(1)}$.
\end{lemma}
\begin{proof}
Use the entropy upper gains
\[
 \alpha_j=t\sqrt{\Gain(D_j)},\qquad
 S_j=\prod_{i\le j}\alpha_i,\quad S_0=1,\qquad \Gamma=S_H,
\]
where $\Gain(D)-1-\log\Gain(D)=2D$. They always use the
\emph{original} entropies, including after replacement.
Every original or replacement row has second-moment operator
norm at most $\alpha_j^2$, so its feature norm is at most
$\sqrt2S_j$. The tangent bound gives
\[
 \log\alpha_j\le(D_j-c_t)/a,\qquad \Gamma\le e^{E/a}.
\]
Write $W=\sqrt{2t^2r_t}$. The margin at any input implies
$W\Gamma\ge1/(2\sqrt2)$, while $W\le\sqrt{2t^2A}$.
The normalized readout multiplier and needed inverse margin
tolerances are therefore polynomial in $R$.

Take $\rho_0=[CH(A+1)]^{-1}\le1/64$. Select common anchors
as in Lemma~\ref{lem:critical-anchors}, with entropy cutoff
$\eta$ and readout projection tolerance $\rho_0$. There are
at most $A/\eta+W^2/\rho_0^2$ anchors. Replace each density
by its conditional expectation $w_j(G_T)$ given the anchors.

A one-dimensional Gaussian $G$ of variance $\sigma^2$ satisfies
\begin{equation}
 W_2^2(\nu,G)\le2\sigma^2\KL(\nu\Vert G).
 \label{eq:prefactor-transport}
\end{equation}
For a smooth increasing transport $T$ from a centered Gaussian,
change of variables and Gaussian integration by parts give
\[
 \KL(\nu\Vert G)=
 \frac{\E(T(X)-X)^2}{2\sigma^2}
 +\E[T'(X)-1-\log T'(X)].
\]
This proves \eqref{eq:prefactor-transport}; translation and
approximation give the general form. Conditional on the anchors,
the reference query variance is at most $t^2Q_{j-1}(x,x)$,
and conditional entropy averaged under the tilted anchor law is
at most $\eta$. Coupling original and replaced queries sharing
that anchor vector gives feature root mean square error at most
$\sqrt{2t^2Q_{j-1}(x,x)\eta}$. After division by $S_j$ this is
at most $2(t/\alpha_j)\sqrt\eta\le2\sqrt\eta$.
Zero conditional variance forces equality and needs no transport.

For transfer to a changed preceding kernel with the same anchor
block, tilt the Gaussian query coupling by the common anchor
density. Its entropy is at most $D_j$ and its root mean square
gain is at most $\alpha_j$. Normalized propagation thus has
gain at most one. Successive replacements define a genuine
full-cube path of normalized query error at most $CH\sqrt\eta$.
Choose $\eta$ so $W\Gamma CH\sqrt\eta\le\rho_0$.
Its inverse and the anchor count are polynomial in $R$.

The anchored readout loses at most $2\rho_0$ margin. Scaling by
$(1-4\rho_0)^{-1}$ restores margin $1/2$, adds
$O(A\rho_0)=O(1/H)$ action, and gives new readout norm
$W'\le2W$. Apply the canonical construction, retaining the span
of the Gaussian coordinate functions and anchor features and
appending the residual norm. It adds at most $|T|+1$ coordinates
per layer. Anchor blocks, query-anchor moments, and diagonals
agree inductively with the genuine replaced path. These are
precisely the invariants used by subsequent anchored densities
and the final readout; full query-query covariance preservation
is not asserted.
\end{proof}

\begin{lemma}[Evaluation after arbitrary scalar normalization]
\label{lem:prefactor-evaluation}
Let $F$ be a canonical map on $\R^K$ for a controlled row law
with second moment $\Sigma$ and $M\le K+q$ basis functions.
For $\alpha\ge\sqrt{\norm\Sigma_{\mathrm{op}}}$ and $S>0$,
the map $F(Sv)/(S\alpha)$ has a constant-threshold-depth,
polynomial-size evaluator on $\norm v\le U$, to any inverse
polynomial accuracy, whenever $K,M,U$ and inverse accuracy are
polynomially bounded. No bound on $S,S^{-1}$ or the density is
needed.
\end{lemma}
\begin{proof}
Let $e(g)\in\R^M$ list an orthonormal basis of the canonical
subspace, and put
\[
 k(g)=\norm g/\alpha,\qquad
 \psi_g(s)=\frac{\tanh(S\norm g\,s)}{S\norm g}.
\]
Assign zero at $g=0$, a null event. Each $\psi_g$ is
zero-preserving and $1$-Lipschitz, hence bounded by $U$ on
$[-U,U]$. The normalized feature is
$k(g)\psi_g(\langle\widehat g,v\rangle)$, with
$\widehat g=g/\norm g$. Use the importance density
\[
 p(g)=\frac{1+\norm{e(g)}^2+k(g)^2}{Z},\qquad
 Z=1+M+\E k^2\le1+M+K.
\]
The projection vector $b(v)$ and squared feature norm $q(v)$
are expectations, under this sampling law, of
\[
 \frac{e(g)k(g)}{p(g)}\psi_g(\langle\widehat g,v\rangle),
 \qquad
 \frac{k(g)^2}{p(g)}\psi_g(\langle\widehat g,v\rangle)^2.
\]
Pointwise $\norm{ek/p}\le Z/2$ and $k^2/p\le Z$.
The vector summand is bounded by $ZU/2$ and has Lipschitz
constant at most $Z/2$; the scalar summand is bounded by
$ZU^2$ and has Lipschitz constant at most $2ZU$.
Bounded concentration and a net of logarithmic cardinality
$O(K\log(\poly(K,M,U)/\eta))$ give polynomially many nodes
approximating both expectations uniformly to tolerance $\eta$.
One may apply scalar concentration to each of the polynomially
many vector coordinates. Complete the state by
\[
 \bigl(b(v),\sqrt{\max\{0,q(v)-\norm{b(v)}^2\}}\bigr).
\]
The square-root inequality only squares the required accuracy.

All sampled directions and coefficients have polynomial range.
Replace each $\psi_g$ by values on a polynomial mesh in
$[-U,U]$. Its Lipschitz bound controls the error independently
of $S\norm g$. The bounded arithmetic and scalar lookup
constructions proved earlier implement these sums, products,
and the clipped square root at constant threshold depth.
The physical scales never become circuit inputs. The
construction is nonuniform.
\end{proof}

\subsection{Structural and actual-posterior converses}

\begin{theorem}[Fixed-prefactor near-baseline parallelization]
\label{thm:prefactor-normal-form}
Fix $0<t<1$. If a full-cube path for $f$ has signed margin
$1/2$ and action $A\le Hc_t+E$, $E\ge0$, then, with
$R=(d+H+A+2)e^E$, the target has threshold size $R^{O_t(1)}$
and depth $O_t(\log R)$. Its Boolean depth is $O_t(\log^2R)$
with polynomial size in $R$.
\end{theorem}
\begin{proof}
Use the anchored canonical path from
Lemma~\ref{lem:prefactor-anchors}, charge each row to its
original $D_j$, and normalize by $S_j$ from that proof.
Each normalized map is globally $1$-Lipschitz and
zero-preserving. Normalized cube states therefore have norm
at most $\sqrt2$ and polynomial dimension.

Let $b_i$ be the two-ball factor of the actual canonical row
covariance. Cover the first $j$ layers, then propagate each
ball through the remaining layers using the entropy upper
gains $\alpha_i$. The possible output error is at most
\begin{equation}
 \sqrt2\,W'\prod_{i\le j}b_i\prod_{i>j}\alpha_i
 \le\sqrt2\,W'
 \exp\left(\frac{A+C/H-Hc_t-j\delta_t}{a}\right).
 \label{eq:prefactor-prefix-error}
\end{equation}
Every layer is charged exactly once. Choose
$J=O_t(E+\log(A+2)+1)$ making this at most $1/64$.
If $J\ge H$, retain all $H$ layers. Otherwise there are at
most $2^J$ prefix balls. Discard empty balls and replace each
remaining center by an actual cube state in it, recording its
label. This doubles the radius and gives stored normalized
centers of norm at most $\sqrt2$. Intermediate existential
cover centers need not be inside the cube-state ball; their
images use the globally defined canonical maps.

The normalized suffix has output Lipschitz constant at most
$W'\Gamma$, polynomial in $R$. Opposite labels at the retained
layer are therefore separated by at least $1/(W'\Gamma)$,
while every query is within $1/(32W'\Gamma)$ of a stored
center. Declare it positive when some positive center is within
$1/(4W'\Gamma)$. Inverse-polynomial state, center and
squared-distance errors preserve this exact classifier.

Evaluate each retained normalized map by
Lemma~\ref{lem:prefactor-evaluation}, projecting radially onto
the radius-$\sqrt2$ ball after each approximate map.
Projection is nonexpansive and fixes the true state, so errors
add across the $O_t(\log R)$ maps. Its arithmetic uses a
reciprocal only above the fixed positive radius threshold and
has constant threshold depth at the required precision.
Distance tests run in parallel. Size is polynomial and
threshold depth logarithmic; the previous Boolean majority
simulation gives squared-logarithmic depth.
\end{proof}

We state the finite-width target-mass input explicitly.
The companion scale-change construction
\cite[Proposition \texttt{prop:long} and the Fourier-seed
section]{uniformtargets} shows that for
$R_g=d+L+A_F(g)+2$, $L\ge C\log R_g$, $N\ge R_g^C$,
and $0<t\le1$, the raw probability of the original strict
global-margin event for $g$ is at least
\begin{equation}
 \exp\left[-N\left\{
 Hc_t+V_*+C\left(a\log R_g+\frac{\log^2R_g}{H}\right)
 \right\}\right].
 \label{eq:prefactor-target-mass}
\end{equation}
Its constants are independent of $t$. This is an explicit
finite-width target construction, not a rate-limit assertion.
As at criticality, target mass is an analytic input; the
extraction and parallelization arguments here are separate.

\begin{theorem}[Fixed-prefactor full-cube posterior ceiling]
\label{thm:prefactor-ceiling}
Fix $0<t<1$, and let $R_0=d+n+L+2$, $d,n\ge1$ and
$L\ge C\log R_0$. Define
\[
 A_0=Hc_t+V_*+
 C\left(a\log R_0+\frac{\log^2R_0}{H}\right),\quad H=L-1.
\]
There are constants $C_t,K_t$ such that at every
\[
 N\ge2^{C_t d}R_0^{C_t}+C_tH\log(1/\beta),\quad0<\beta<1,
\]
the original posterior at $\sigma^2=t^2/N$ gives mass at least
$1-\beta$, simultaneously for every consistent sample of at
most $n$ points, to the sample-independent class
\[
 \mathcal G_t=\{g:J_{t,L}(g)\le A_0+K_t/H\}.
\]
Every member has threshold size $R_0^{O_t(1)}$ and depth
$O_t(\log R_0)$.
\end{theorem}
\begin{proof}
The sample memorizer from
Theorem~\ref{thm:critical-ceiling} has Fourier variation at
most $1+2n$. Equation~\eqref{eq:prefactor-target-mass} bounds
the conditioning denominator below by $e^{-NA_0}$.

Lemma~\ref{lem:critical-dual-certificate} applies unchanged
with references $\mathcal N(0,t^2Q)$ and readout cost $r_t$.
Hidden observables remain bounded; given them the readout has
covariance $t^2Q_H/N$ and
$\E e^{(1-s)Nr_t}\le s^{-m/2}$. Its width formula is unchanged.
Consistent-path transfer now has recurrence
$e_j\le\alpha_j e_{j-1}+\sqrt{m\Delta}$ and final multiplier
$\sqrt{2t^2r_t}$. For total cost $B=O_t(H+\log R_0)$,
each $\alpha_j=t\sqrt{\Gain(D_j)}$ is polynomial in $B+1$.
Suffix products may be larger, but their logarithms are at
most $O_t(H\log(B+2))$. Hence a tolerance with
\[
 \log(1/\Delta)\le
 C_t\{d+H\log(B+2)+\log(H(B+2))\}
\]
makes the output error at most
$\rho=[C_tH(B+1)]^{-1}$. The embedding and data-processing
proof of Lemma~\ref{lem:critical-consistent-transfer}
preserves the action. Readout scaling by $(1-2\rho)^{-1}$
restores margin $1/2$ and adds $O_t(1/H)$ action.

Consider the bad raw event of globally margined draws with
rate above $A_0+K_t/H$. If its probability were at least
$e^{-N(A_0+1/H)}$, use the full-cube certificate with
$m=2^d$, $\tau=1/H$ and the preceding tolerance.
Its width requirement is at most $2^{C_t d}R_0^{C_t}$.
The repaired path would give a bad function action at most
$A_0+K'_t/H$, contradicting $K_t>K'_t$.
The bad raw mass is therefore less than $e^{-N(A_0+1/H)}$.
Division by the denominator gives posterior bad mass at most
$e^{-N/H}\le\beta$.

Since $A_0-Hc_t=O_t(\log R_0)$ under the depth condition,
Theorem~\ref{thm:prefactor-normal-form}, applied within $1/H$
of each rate infimum, supplies uniform circuit bounds.
No attainment assumption is needed.
\end{proof}

\begin{corollary}[Fixed-prefactor sample-to-circuit converse]
\label{cor:prefactor-converse}
Fix $0<t<1$, $0\le\varepsilon<1/2$, and
$\kappa=1-2\varepsilon$. Let $L\ge C\log R_0$,
$R_0=d+n+L+2$, and fix one width above the preceding
ceiling with $\beta=1/6$. If every input law gives positive
probability to an $n$-sample with confidence-failure mass at
most $\varepsilon$, at $\sigma^2=t^2/N$, then $f$ has
threshold size polynomial in $R_0+\kappa^{-1}$ and depth
$O_t(\log R_0)$. Its Boolean depth is
$O_t(\log^2R_0+\log(\kappa^{-1}+2))$, with polynomial size.
\end{corollary}
\begin{proof}
Use the two-majority and minimax proof of
Theorem~\ref{thm:converse} with the fixed class $\mathcal G_t$.
The width is chosen before the input law, the ceiling is
simultaneous over all samples, and
Theorem~\ref{thm:prefactor-normal-form} supplies the uniform
voter circuit bounds.
\end{proof}

\begin{corollary}[Uniformity away from the critical prefactor]
\label{cor:prefactor-uniform}
Fix $0<t_0<1$. The structural theorem holds uniformly for
$0<t\le t_0$ with resource
\[
 R=(d+H+t^2A+2)\exp(E/a).
\]
The posterior ceiling and sample converse hold uniformly on this
same range: their polynomial exponents and width constants depend
only on $t_0$, if the ceiling class is replaced by
\[
 \mathcal G_t^{\,\mathrm{unif}}
 =\{g:J_{t,L}(g)\le A_0+K/(t^2H)\}.
\]
In particular the finite-width threshold can be chosen as
$2^{C_{t_0}d}R_0^{C_{t_0}}+C_{t_0}H\log(1/\beta)$,
independently of $t$.
\end{corollary}
\begin{proof}
Put $Q=t^2A$ and retain the entropy upper gains from
Lemma~\ref{lem:prefactor-anchors}. Use layer-dependent entropy
cutoffs $\eta_j=\eta\Gain(D_j)$. Their total anchor count is
at most
\[
 \sum_j\frac{D_j}{\eta\Gain(D_j)}
 \le\frac{H}{2\eta},
\]
because $D/\Gain(D)\le1/2$. The normalized conditional
transport error is now
$2(t/\alpha_j)\sqrt{\eta_j}=2\sqrt\eta$.
Take $\rho_0=c/[H(1+Q)]$. Since $W^2\le2Q$ and
$\Gamma\le e^{E/a}$, the required $\eta^{-1}$ and readout
projection count are polynomial in the displayed $R$,
uniformly in $t$. Readout repair adds at most
$CA\rho_0\le C/(t^2H)$ action and at most doubles its norm.
The cover still charges the original hidden costs, so this
norm factor is the only change needed in
\eqref{eq:prefactor-prefix-error}.

As a function of $a$, $\delta$ has $\delta(0)=0$ and
$\delta''(a)=1/[2(a+1)(a+2)]>0$. Thus $\delta(a)/a$ is
increasing and bounded below by its positive value at
$a_0=t_0^{-2}-1$. The mixed cover error is bounded by
$\sqrt2W'\exp(E/a-j\delta_t/a)$, and $O_{t_0}(\log R)$
layers suffice. The arbitrary-scale evaluator requires no
additional bound involving $t^{-1}$.

For the statistical assertion, the target budget satisfies
\[
 t^2A_0\le C(H+\log R_0+1),\qquad
 \frac{A_0-Hc_t}{a}\le C_{t_0}\log R_0.
\]
These bounds persist after addition of $K/(t^2H)$, since
$(t^2a)^{-1}=(1+a)/a\le1+a_0^{-1}$.
In the finite-event certificate take
$\tau=1/(t^2H)$ and event budget $B=A_0+\tau$.
Then $B/\tau=Ht^2B$ is polynomial in $R_0$ and
$1/\tau=t^2H\le H$. Transfer gains satisfy
$\alpha_j^2\le4t^2B+2$, and the readout norm is at most
$\sqrt{2t^2B}$. With repair tolerance
$\rho=c/[H(1+t^2B)]$, the required
$\log(1/\Delta)$ has a polynomial bound independent of $t$.
The certificate's width bound is therefore uniform, and
readout repair adds $O(1/(t^2H))$ action.

The bad-event contradiction now uses raw threshold
$\exp[-N(A_0+1/(t^2H))]$ and bad rate
$A_0+K/(t^2H)$. Its posterior mass is at most
$e^{-N/(t^2H)}\le e^{-N/H}$. The structural resource of
every retained function is polynomial in $R_0$, with constants
depending only on $t_0$. The same two-majority argument proves
the uniform sample converse.
\end{proof}

Thus arbitrarily small prefactors do not cause a loss of uniformity
on a range bounded away from one. A prescribed prefactor may vary
with dimension or width within that range. The direct spectral
proof deteriorates at $t\to1$; the next section resolves that
boundary by combining block covering with critical parallelization.

\section{Uniform parallelization across the critical boundary}
\label{sec:uniform-prefactor}

We now remove the remaining dependence on a prefactor approaching
one. The spectral argument alone would retain $O(S/a)$ layers
and could create $2^{O(S/a)}$ representatives. Instead we cover
blocks of $O(1/a)$ physical layers with a constant number of
balls. A low-entropy prefix then has only $O(S)$ such blocks,
and the critical parallel iteration evaluates the whole prefix.

\begin{theorem}[Uniform near-baseline parallelization]
\label{thm:uniform-normal-form}
Fix a numerical constant $K_0$. There are constants $C_0,C$,
depending only on $K_0$, such that the following holds for
every $0<t\le1$. Suppose a full-cube path for $f$ has signed
margin $1/2$ and
\[
 A\le Hc_t+V_*+K_0\left(aS+\frac{S^2}{H}\right),
 \quad S\ge1+\log(d+H+2),\quad H\ge C_0S,
\]
where $a=t^{-2}-1$ and $c_t=(a-\log(1+a))/2$.
Then $f$ has threshold size $\exp(CS)$ and depth $CS$.
Its Boolean size is $\exp(CS)$ and depth at most $CS^2$.
All constants are independent of $t$.
\end{theorem}

\subsection{Adjusted costs and a controlled prefix}

We first supply the scalar estimates used in its proof. The
critical value $V$ is nonincreasing: starting from a larger
variance, use the same standardized law at each scalar step.
This preserves its entropy cost and increases every subsequent
variance by monotonicity of $\tanh^2$ under dilation. The
terminal readout cost decreases. Taking infima proves the claim.

Put $q_j=\max_xQ_j(x,x)$, $q_0=(d+1)/d$, and
\[
 e_j=D_j-c_t-\frac a2\log(q_j/q_{j-1}),\qquad
 r_1=t^2r_t.
\]
All query variances are positive, since zero variance would
propagate at finite entropy and preclude a nonzero final output.
At a query attaining $q_j$, let its preceding variance be
$p\le q_{j-1}$ and put $M=\E G_x^2/p$. The scalar reference
identity is
\[
 D_t(\eta)=D_1(\eta)+c_t+\frac a2(M-1).
\]
Since $q_j\le Mp$, the difference between $e_j$ and this
query's critical scalar entropy is nonnegative. Monotonicity of
$V$ and of $p-cp^2$ therefore gives
\begin{equation}
 e_j\ge0,\quad
 e_j\ge V(q_{j-1})-V(q_j),\quad
 e_j\ge q_j-q_{j-1}+cq_{j-1}^2.
 \label{eq:uniform-adjusted}
\end{equation}
Moreover
\[
 \widetilde A:=\sum_j e_j+r_1
 =A-Hc_t-\frac a2\log(q_H/q_0)-ar_1
 \le A-Hc_t+a.
\]
The last inequality follows from
$r_1\ge1/(8q_H)$ and
$\tfrac12\log(q_H/2)+1/(8q_H)\ge-1$.
Consequently, for $\epsilon=C_{K_0}(aS+S^2/H)$,
\begin{equation}
 \widetilde A\le V_*+\epsilon,\qquad
 \sum_{j\le i}e_j\le\epsilon+Cq_i.
 \label{eq:uniform-prefix-budget}
\end{equation}
The latter follows by telescoping the Bellman inequalities over
the suffix and using $r_1\ge V(q_H)$.

Use the original entropy upper gains
\[
 \alpha_j=t\sqrt{\Gain(D_j)},\quad
 P_j=\prod_{i\le j}\alpha_i,\quad P_0=1,\quad
 \Gamma=P_H,\quad g_j=\log\alpha_j.
\]
Their exact inverse relation gives
\begin{equation}
 D_j-c_t=ag_j+
       \frac{1+a}{2}(e^{2g_j}-1-2g_j),\qquad
 0\le h_j:=D_j-c_t-ag_j\le e_j.
 \label{eq:uniform-gain-remainder}
\end{equation}
Here $q_j/q_{j-1}\le\alpha_j^2$ proves the last inequality.

\begin{lemma}[Uniform gain and last-exit bounds]
\label{lem:uniform-gain-prefix}
Under the hypotheses of Theorem~\ref{thm:uniform-normal-form},
$\Gamma\le e^{CS}$ and $t^2A\le CH$. For $0<a\le1/4$,
the last index $\tau$ with $q_\tau\le a$ satisfies
\[
 H-\tau\le C(\epsilon/a^2+1/a).
\]
If no such index exists, the same bound holds for $H$.
\end{lemma}
\begin{proof}
Write $u_j=\log(q_j/q_{j-1})$. When $u_j>0$,
the critical entropy dominated by $e_j$ and Gaussian
second-moment duality imply $e_j\ge u_j^2/4$.
Thus $u_j\le2\sqrt{e_j}$. If $g_j>\sqrt{e_j}\ge0$,
the nonnegative remainder in
\eqref{eq:uniform-gain-remainder} is at least
$(1+a)g_j^2$, contradicting
$D_j-c_t=au_j/2+e_j$.
Hence $g_j\le\sqrt{e_j}$. The critical occupation proof,
using \eqref{eq:uniform-adjusted}--\eqref{eq:uniform-prefix-budget},
gives
\[
 \sum_jg_j\le C\log(H+2)+C\sqrt{H\epsilon}.
 \tag{*}
\]
There is also the telescoping tangent inequality
$g_j\le u_j/2+e_j/a$.

For clarity, its complementary last-exit estimate is
\[
 \sum_jg_j\le C\{1+\log(1/a)+\epsilon/a\},
 \qquad 0<a<1/4,\quad\epsilon<1/4.
 \tag{**}
\]
To verify it, put $m=\min_iq_i$ and start dyadic thresholds
at $b=\max(a,m)$. If $m\le a$, the gain through the last
state below $a$ is at most $C+\epsilon/a$ by the tangent
inequality and the prefix budget. If $m>a$, the reciprocal
occupation bound proved in the critical section gives initial
length at most $C(\epsilon/m^2+1/m)$ and cost at most
$\epsilon+Cm$; Cauchy--Schwarz bounds its gain by
$C(1+\epsilon/m)$. Between the last exits below $u$ and
$2u$, remove one initial transition. Pre-states exceed $u$,
the final state is at most $2u$, and the cost is at most
$\epsilon+2Cu$. Reciprocal occupation and Cauchy--Schwarz
give length $C(\epsilon/u^2+1/u)$ and gain
$C(1+\epsilon/u)$. Sum geometrically up to a fixed positive
threshold; the remaining tail has bounded length and gain.
The same occupation calculation starting at threshold $a$
proves the asserted last-exit length, without requiring small
$\epsilon$. If $m>a$, include the initial minimum-state
prefix. This also covers the case with no state below $a$.

For the required gain, if $a\le S/H$, use (*) to obtain
$CS$. Otherwise $\epsilon/a\le CS$ and $a\ge1/H$.
If $\epsilon\ge1/4$ or $a\ge1/4$, the whole-path tangent
bound $(V_*+\epsilon)/a$ is $O(1+\epsilon/a)$ in the
respective case. In the remaining case use (**), with
$\log(1/a)\le\log H\le S$. This proves the uniform gain.
Finally $t^2c_t\le1/2$, $t^2a\le1$, and $S\le H/C_0$
give $t^2A\le CH$.
\end{proof}

We prove Theorem~\ref{thm:uniform-normal-form}. At $t=1$ use
Theorem~\ref{thm:critical-normal-form}; at $a\ge1/4$ use
Corollary~\ref{cor:prefactor-uniform}, whose resource has
logarithm $O(S)$. Hence suppose $0<a<1/4$.

Choose a large fixed $K$ below and set $J=\lceil KS/a\rceil$.
If $H\ge K_HS/a$, where $K_H$ is sufficiently large compared
with $K$ and the last-exit constants, then $J\le\tau$ and
\begin{align}
 E_{\rm pre}:=\sum_{j\le J}e_j&\le C_eaS,\nonumber\\
 D_{\rm pre}:=\sum_{j\le J}D_j
 &\le Jc_t+C_eaS=O_K(aS),&
 \sum_{j\le J}\sqrt{D_j}&=O_K(S).
 \label{eq:uniform-prefix-cost}
\end{align}
The constant $C_e$ is independent of $K$, since
$S^2/H\le aS/K_H$ and $q_J$ is controlled by the last-exit
prefix budget. More precisely, nonnegativity gives
$E_{\rm pre}\le\sum_{j\le\tau}e_j\le\epsilon+Ca$.
Also $q_J\le1\le q_0$, so the telescoping logarithm in
$D_{\rm pre}$ is nonpositive.

We need the sharper variance estimate
\begin{equation}
 \sum_{j<J}q_j\le C\log(J+2)+C\sqrt{JE_{\rm pre}}
                  \le C(1+\sqrt K)S.
 \label{eq:uniform-prefix-occupation}
\end{equation}
The forced-cooling comparison from the critical proof works for
arbitrary total forcing by taking
$r_{j+1}=r_j-cr_j^2$ and
$s_{j+1}=\min\{1,s_j-cs_j^2+e_{j+1}\}$.
On a capped step $q_{j+1}\le1\le r_{j+1}+s_{j+1}$.
On an uncapped step, monotonicity and subadditivity of
$u-cu^2$ on $[0,3]$ preserve $q_j\le r_j+s_j$.
Capping preserves $c\sum s_j^2\le\sum e_j$.
Now $\sum r_j=O(\log(J+2))$ and Cauchy--Schwarz give
\eqref{eq:uniform-prefix-occupation}.

If instead $H<K_HS/a$, then
$A\le V_*+C_{K_H}S^2/H$.
The full critical canonical/Picard proof applies to the raw
subcritical costs with uniform constants. Indeed its scalar
Bellman comparison is trivial when the next variance is no
larger; when it is larger, $M>1$ in the reference identity
and $D_t\ge D_1$. Its raw dissipation is
\[
 D_t\ge q-t^2p+ct^4p^2\ge q-p+c'p^2
\]
since $t^2\ge4/5$. Gaussian gains are at most the critical
gains, and the residual estimate after substituting $g=t\xi$
is at most $C(p+\sqrt{D_t})$. The readout cost also dominates
the critical cost. Thus the occupation, anchor and Picard
proof gives size $\exp(CS)$ and depth $CS$ in this short
regime. It remains to treat the long prefix above.

\subsection{Gaussian and non-Gaussian row costs}

Write $W=\sqrt{2t^2r_t}$. Besides $W\Gamma\le e^{CS}$,
the absolute margin bounds every partial normalization:
\begin{equation}
 P_j\le2W\Gamma\le e^{CS}\qquad(j\ge1).
 \label{eq:uniform-partial-gain}
\end{equation}
Indeed $q_H(x)\le q_j(x)(\Gamma/P_j)^2$,
$q_j(x)\le1$, and $1/2\le W\sqrt{q_H(x)}$.
Perform the adaptive normalized anchor replacement from
Corollary~\ref{cor:prefactor-uniform}, using the original
$\alpha_j$. Take readout tolerance
$\rho=c/[H(1+t^2A)]$ and choose the entropy cutoff parameter
$\eta$ small enough that
\[
 W\Gamma CH\sqrt\eta\le\rho,\qquad
 CH\sqrt\eta\max_jP_j\le H^{-2}.
\]
The inverse cutoff, anchor count, and canonical dimensions are
all $\exp(O(S))$. The genuine replaced path has maximum
diagonals $\widehat q_j$ satisfying
\begin{equation}
 \sum_{j<J}\widehat q_j\le\sum_{j<J}q_j+C.
 \label{eq:uniform-variance-transfer}
\end{equation}
This follows from
$|\sqrt{\widehat q_j}-\sqrt{q_j}|\le H^{-2}$ and bounded
hidden variances. Its repaired readout has norm $W'\le2W$.
The added action $O(1/(t^2H))$ is absorbed by the theorem's
budget after increasing $K_0$.

For any original controlled Gaussian row, decompose its entropy
using its full second moment $\Sigma$:
\[
 D=J_{\rm ng}+C_{\rm g},\qquad
 J_{\rm ng}=\KL(\rho\Vert\mathcal N(0,\Sigma)),\qquad
 C_{\rm g}=\sum_iF_t(\lambda_i(\Sigma)).
\]
At the query attaining $q_j$, the Gaussian cost is at least
$F_t(M)$, with $q_j/q_{j-1}\le M$.
The tangent inequality for $F_t$ implies
$J_{{\rm ng},j}^{\rm original}\le e_j$.
Anchor replacement does not increase this non-Gaussian entropy.
To see this, split Gaussian coordinates orthogonally into the
anchor span and its complement. The replaced law is its
controlled anchor marginal times an independent free Gaussian.
Its matching Gaussian has the same product form, so its
$J_{\rm ng}$ equals that of the anchor marginal.
Data processing of the original law and its matching Gaussian
proves the inequality. Changed kernels with exact anchor block
alter only the isometric realization and the free complement.
The conclusion therefore persists in canonical coordinates.
For the canonical row costs we have
\begin{equation}
 J_{{\rm ng},j}+C_{{\rm g},j}\le D_j,\qquad
 \sum_{j\le J}J_{{\rm ng},j}\le C_eaS.
 \label{eq:uniform-ng-budget}
\end{equation}

\subsection{Blocks with a constant number of representatives}

For a physical canonical map $F_j$, let $A_j$ be its linear
part, obtained from the retained Gaussian coordinate functions.
Then $A_j^TA_j=\Sigma_j$. Its polar decomposition gives
\[
 A_j=tU_j(I+B_j),\qquad
 \norm{B_j}_{\rm HS}^2\le2C_{{\rm g},j},
\]
where $U_j$ is an isometry. The last inequality is
$(\sqrt\mu-1)^2\le\mu-1-\log\mu$ summed over eigenvalues.
For a block of $b$ layers, telescoping yields
\begin{equation}
 A_{\rm block}=t^bU_{\rm block}+E_{\rm block},\qquad
 \norm{E_{\rm block}}_{\rm HS}
 \le e^{\sqrt{2bC_{\rm g,block}}}
                   \sqrt{2bC_{\rm g,block}}.
 \label{eq:uniform-block-linear}
\end{equation}
This holds also for increasing rectangular dimensions, since
products of the $U_j$ remain isometries.

The nonlinear residual $N_j=F_j-A_j$, on its physical state
ball of squared radius $\widehat q_{j-1}$, satisfies
\begin{equation}
 \Lip(N_j)\le C\sqrt{\lambda_j}
       \bigl(\lambda_j\widehat q_{j-1}+
                         \sqrt{J_{{\rm ng},j}}\bigr),
 \qquad\lambda_j=\norm{\Sigma_j}_{\rm op}.
 \label{eq:uniform-matched-residual}
\end{equation}
Whiten by the matching Gaussian, $g=\Sigma_j^{1/2}\xi$.
The law of $\xi$ has entropy $J_{{\rm ng},j}$ relative to
a standard Gaussian. The residual moment estimate in
Lemma~\ref{lem:critical-canonical}, followed by composition
with $\Sigma_j^{1/2}$, proves
\eqref{eq:uniform-matched-residual}. The canonical residual
projection is a Hilbert-space contraction.

Set $b=\lceil c/a\rceil$, with a large fixed $c$, and call a
full block Gaussian-good if
$C_{\rm g,block}\le r_0a$. Then
$bC_{\rm g,block}\le(c+1/4)r_0$, so
\eqref{eq:uniform-block-linear} has a fixed bound $M_0$.
Every subproduct of the actual linear gains has a fixed
bound too, since
$\sum_{\rm block}\sqrt{C_{{\rm g},j}}
\le\sqrt{bC_{\rm g,block}}$.
Telescoping $F_{\rm block}-A_{\rm block}$ on pairs of
\emph{actual query states} gives
\[
 \Lip(N_{\rm block})\le C_{\rm block}
       \sum_{\rm block}
       (\widehat q_{j-1}+\sqrt{J_{{\rm ng},j}}).
 \tag{\(\dagger\)}
\]
All intermediate states for these pairs lie in the balls
used in \eqref{eq:uniform-matched-residual}; both the full
and linear maps have gain at most $\sqrt{\lambda_j}$.
No ambient-ball assertion about this residual is needed.

Choose $c$ so large that
\[
 \sigma_0=e^{-c/(2(1+1/4))}
   < B_0/32,\qquad B_0=e^{-(c+1/4)/4-1}.
\]
Since $c_t\le a^2/4$, for every block cost
$D_{\rm block}=\sum_{\rm block}D_j$,
\[
 t^b\le\sigma_0,\qquad
 \exp((D_{\rm block}-bc_t)/a-1)\ge B_0.
\]
In \eqref{eq:uniform-block-linear}, truncate singular values
of $E_{\rm block}$ below a sufficiently small fixed multiple
of $B_0$. The retained rank is bounded by a constant depending
only on $c,r_0$. A volumetric net for that fixed-rank ellipsoid
therefore covers the image of any radius-$r$ ball under
$A_{\rm block}$ by at most $M_{\rm cover}$ balls of radius
$B_0r/8$, where $M_{\rm cover}$ is an absolute constant.
The free isometry contributes at most $\sigma_0r$,
and the discarded operator and net errors each contribute
a sufficiently small fixed multiple of $B_0r$.

Also call a block nonlinear-good if
\[
 \sum_{\rm block}
 (\widehat q_{j-1}+\sqrt{J_{{\rm ng},j}})\le\eta_{\rm nl},
\]
where $\eta_{\rm nl}>0$ is fixed so that
(\(\dagger\)) is at most $B_0/16$. In each nonempty linear
cover cell select an actual query state. Two such states
within the original cluster have distance at most $2r$,
their linear images differ by at most $B_0r/4$, and their
nonlinear remainders differ by at most $B_0r/8$.
Thus the true outputs have an actual-representative cover of
radius at most
\begin{equation}
 \frac{B_0r}{2}
 \le r\exp((D_{\rm block}-bc_t)/a-1),
 \label{eq:uniform-good-block}
\end{equation}
with at most $M_{\rm cover}$ cells. The selection of actual
representatives is what permits the restricted estimate
(\(\dagger\)). On any bad block, simply propagate the
existing center and use the global one-ball gain
$\exp((D_{\rm block}-bc_t)/a)$.
Treat the final incomplete block the same way.

There are $M=\lfloor J/b\rfloor$ full blocks, with
$M\ge KS/[2(c+1/4)]$ after increasing $K$.
By \eqref{eq:uniform-prefix-cost}, at most
$(K/4+C_e+1)S/r_0$ are Gaussian-bad.
Choose $r_0$ sufficiently large compared with $c$.
Equations \eqref{eq:uniform-prefix-occupation},
\eqref{eq:uniform-variance-transfer} and
\eqref{eq:uniform-ng-budget} give
\[
 \sum_{j\le J}
 (\widehat q_{j-1}+\sqrt{J_{{\rm ng},j}})
 \le C(1+\sqrt K)S.
\]
After $c,r_0,\eta_{\rm nl}$ are fixed, choose $K$ large
enough that each type of bad block accounts for at most
$M/4$. At least $KS/[4(c+1/4)]$ blocks are good.
Finally choose $K_H$ to validate the long-prefix selection.
This order of choices involves no problem parameter.

Start with an actual cube state and radius $2\sqrt2$.
Multiplying good and bad block gains and dividing the prefix
state by $P_J$ gives normalized cover radius at most
\[
 2\sqrt2\exp\left(\frac{\sum_{j\le J}h_j}{a}
                      -G_{\rm good}\right)
 \le2\sqrt2\exp(C_eS-G_{\rm good}).
\]
Enlarge fixed $K$ so the subsequent readout multiplier
$W'\Gamma\le e^{CS}$ leaves output uncertainty at most
$1/64$. The number of stored representatives is
$M_{\rm cover}^M=\exp(O(S))$, since $M=O(S)$.
Every stored normalized center is an actual cube state,
of norm at most $\sqrt2$. The normalized suffix has output
Lipschitz constant at most $W'\Gamma$. Opposite labels
are separated by $1/(W'\Gamma)$, so the same distance
classifier as before is exact at inverse-exponential-in-$S$
coordinate accuracy.

\subsection{Parallel evaluation and the finite-width consequence}

The prefix may contain $J\gg S$ layers, but
\eqref{eq:uniform-prefix-cost} and
\eqref{eq:uniform-prefix-occupation} give
\[
 \sum_{j\le J}\sqrt{D_j}+
 \sum_{j<J}\widehat q_j=O_K(S).
\]
The raw-reference residual estimate is
$\Lip(F_j-A_j)\le C(\widehat q_{j-1}+\sqrt{D_j})$.
After normalization by $P_j$, divide by
$\alpha_j\ge t\ge\sqrt{4/5}$; its total Lipschitz budget
is still $O(S)$ and the linear maps are contractions.
Extend each residual by radial projection onto its actual
normalized state ball. Floor that radius at $e^{-C_1S}$.
By \eqref{eq:uniform-partial-gain}, the total increase in
physical squared radii is at most
$H e^{CS}e^{-2C_1S}$, made at most one by fixed large $C_1$.
This adds only a constant to the residual budget and leaves
all actual states unchanged.

Precompute all interval products of the linear contractions
and apply the strict-causal Picard iteration from
Theorem~\ref{thm:critical-normal-form}. Its remainder is
bounded by $C\Lambda^k/k!$, with $\Lambda=O(S)$, and
uniform evaluation errors are amplified by at most
$e^\Lambda$. Thus $O(S)$ iterations suffice at
$e^{-CS}$ primitive accuracy.
All time indices and state dimensions are polynomial in
$e^S$. Lemma~\ref{lem:prefactor-evaluation} evaluates each
normalized canonical map at constant threshold depth,
independently of small $P_j$; subtracting its bounded linear
part and projecting supplies the residual evaluator.
Appending the parallel distance tests proves
Theorem~\ref{thm:uniform-normal-form}.\hfill$\square$

\begin{theorem}[Uniform actual-posterior circuit ceiling]
\label{thm:uniform-ceiling}
There are universal positive constants
\[
 C_{\mathrm{mass}},\quad C_{\mathrm{pad}},\quad C_{\mathrm{width}},
 \quad C_{\mathrm{circ}},\quad K
\]
with the following property.
Put $R_0=d+n+L+2$ and suppose
$L\ge C_{\mathrm{pad}}\log R_0$.
For every $0<t\le1$, define
\[
 A_0=Hc_t+V_*+
 C_{\mathrm{mass}}\left(a\log R_0+\frac{\log^2R_0}{H}\right).
\]
At every actual width
\[
 N\ge2^{C_{\mathrm{width}}d}R_0^{C_{\mathrm{width}}}
       +C_{\mathrm{width}}H\log(1/\beta),\qquad0<\beta<1,
\]
and every consistent sample of at most $n$ points, the
original posterior at $\sigma^2=t^2/N$ gives mass at least
$1-\beta$ to
\[
 \{g:J_{t,L}(g)\le A_0+K/(t^2H)\}.
\]
Every such function has threshold size at most $R_0^{C_{\mathrm{circ}}}$
and depth at most $C_{\mathrm{circ}}\log R_0$.
The mass-budget constant is fixed before choosing the padding, width,
and circuit-resource constants; increasing a resource constant does
not enlarge the defining rate class.
Thus one common class of circuits with those resources is
a posterior ceiling class for all $t,N$ in the assertion.
\end{theorem}
\begin{proof}
Use the uniform target mass
\eqref{eq:prefactor-target-mass} for the sample memorizer.
The finite-event argument of
Corollary~\ref{cor:prefactor-uniform} was uniform in
extraction, without requiring separation from $t=1$:
take normalized action $Q=t^2A$, slack
$\tau=1/(t^2H)$, and repair tolerance
$c/[H(1+Q)]$. The certificate ratios
$A/\tau=HQ$ and $1/\tau=t^2H\le H$, transfer gains
$\alpha_j^2\le4Q+2$, readout norm, and
$\log(1/\Delta)$ all have uniform polynomial bounds.
The raw-event contradiction gives exceptional posterior
mass at most $e^{-N/(t^2H)}\le e^{-N/H}$.
The repaired rate threshold differs from $A_0$ by
$O(1/(t^2H))$, which is absorbed by
$C(aS+S^2/H)$ for $S=C'\log R_0$.
Theorem~\ref{thm:uniform-normal-form} now gives uniform
circuit resources, including for paths approaching a rate
infimum. The class of all circuits with these resource
bounds is independent of $t$ and of the sample.
\end{proof}

\begin{corollary}[Uniform sample-to-parallel-circuit converse]
\label{cor:uniform-converse}
Let $0\le\varepsilon<1/2$, $\kappa=1-2\varepsilon$,
$R_0=d+n+L+2$, and assume the padding hypothesis of
Theorem~\ref{thm:uniform-ceiling}.
Fix one width above that theorem's
threshold with $\beta=1/6$, and any $0<t\le1$.
If every input law has positive probability of an
$n$-sample with confidence-failure mass at most
$\varepsilon$, then $f$ has threshold size polynomial
in $R_0+\kappa^{-1}$ and depth $O(\log R_0)$.
Its Boolean depth is
$O(\log^2R_0+\log(\kappa^{-1}+2))$, with polynomial size.
The constants do not depend on $t$.
\end{corollary}
\begin{proof}
Apply the two-majority and minimax proof of
Theorem~\ref{thm:converse} to the common circuit class
from Theorem~\ref{thm:uniform-ceiling}.
Choose the one width and its actual prefactor before
the input law. The ceiling is simultaneous over samples
and its threshold is independent of the prefactor.
\end{proof}

\section{A universal polynomial bound would imply a collapse}

\begin{lemma}[Bounded-weight tanh teachers for Boolean circuits]
\label{lem:boolean-teacher}
Every bounded-fan-in Boolean circuit of size $s$ on $d$ inputs has a
tanh representation with global signed output margin greater than
one, depth and width polynomial in $d+s$, and every physical weight
an integer of absolute value at most $20$. It has no biases after its augmented
input.
\end{lemma}
\begin{proof}
Use $\{-1,1\}$ wire values. AND, OR, NOT, and COPY have respective
affine forms
\[
 u+v-1,\qquad u+v+1,\qquad -u,\qquad u,
\]
of absolute margin at least one on Boolean inputs. Process gates in
topological order, retaining all earlier values by copy gates; this
gives width $O(d+s)$ and depth $O(d+s)$, with one initial layer if
needed to encode input wires. Replace each gate by tanh of ten times
its affine form. If every incoming value is within $1/100$ of its
intended sign, the affine perturbation is at most $3/100$, and the
new sign error is at most $2e^{-19.4}<1/100$. Thus errors do not
accumulate through the depth.

Implement each affine constant with a separate wire of intended
value $+1$, initialized as $\tanh(10)$ from the augmented constant
input and propagated by the same COPY map. Include its error in
the preceding $3/100$ perturbation bound. Its gate coefficient is
$\pm10$. Duplicate logical wire copies when needed so repeated
gate inputs use separate coordinates; this changes width by only
a constant factor. All hidden weights are then integers of
absolute value at most ten. This also handles constant Boolean
circuits. Finish with twice the encoded output as a linear
readout. Its signed margin exceeds $1.98$. Pad with copy layers if
necessary to obtain weighted depth at least three.

For every larger common width $N$, pad all hidden layers by zero
neurons and use zero readout weights on them. Since $\tanh(0)=0$,
this preserves the classifier, depth, and margin. On the finite
cube, continuity gives an open parameter neighborhood with the
same strict margin. Every nondegenerate Gaussian prior assigns
that neighborhood positive probability, so these widths are
realizable at every finite positive scale.
\end{proof}

\begin{corollary}[Joint polynomial learning implies $\Ppoly=\TC^1$]
\label{cor:collapse}
Fix any prescribed prefactor schedule $0<t=t(d,L,N)\le1$,
chosen independently of the input law and the observed sample.
Suppose there is one polynomial sample bound jointly
in $d,C,L,\varepsilon^{-1},\log(2/\delta)$ with the following property.
Every Boolean target admitting a depth-$L$, width-at-most-$C$ tanh
teacher with weights bounded by $C$ and global signed margin greater
than $1/2$ is learned by the original posterior at scale
$\sigma^2=t(d,L,N)^2/N$, for every input law and every integer
$N\ge C$. The guarantee is that confidence-failure mass is at most
$\varepsilon$ with sample probability at least $1-\delta$.
Then
\[
 \boxed{\qquad\Ppoly=\TC^1.\qquad}
\]
In particular $\Ppoly=\NC^2$. The same conclusion holds if the
hypothesis instead covers every realizable width.
In particular, the conclusion applies to every fixed exponent
$\gamma\ge1$, every fixed prefactor, and arbitrary prescribed
prefactors approaching one.
\end{corollary}
\begin{proof}
Take any family in $\Ppoly$ and a polynomial-size bounded-fan-in
Boolean circuit family for it. Lemma~\ref{lem:boolean-teacher}
supplies teacher width, weight bound, and depth all polynomial in $d$.
Enlarge $C$ to bound both the width and twenty, and pad the teacher
depth to at least $d$ using the same robust copy gates. Fix, for example,
$\varepsilon=1/3$ and $\delta=1/2$. The hypothesized sample count is
$n=\poly(d)$, determined without choosing the learner width.

The padding $L\ge d$ and the polynomial bounds on
$n,L$ ensure $L\ge C_0\log(d+n+L+2)$ for all sufficiently large
$d$. Choose one width above both $C$ and the uniform ceiling
threshold in Theorem~\ref{thm:uniform-ceiling} with $\beta=1/6$.
This threshold is independent of $t$, so the actual value
$t(d,L,N)$ can now be used without circularly changing $N$.
The width and scale have been fixed before the input law.
Corollary~\ref{cor:uniform-converse} gives threshold size
polynomial in $d$ and depth $O(\log d)$, with uniform constants.
No computability assumption on the prescribed schedule is
needed for this nonuniform circuit conclusion.
Finitely many remaining dimensions can be
hardcoded. Hence $\Ppoly\subseteq\TC^1$. The reverse
inclusion holds by definition. The Boolean simulation from
Theorem~\ref{thm:main} also gives
$\TC^1\subseteq\NC^2\subseteq\Ppoly$.
\end{proof}

The width quantifier is essential here. A sample bound polynomial in
the learner width $N$ is not the joint bound in this corollary: it
would make $n$ depend on the width subsequently chosen above the
ceiling threshold. Likewise, the result does not say that a fixed
input distribution admitting polynomial sample complexity forces an
exact parallel circuit.

\begin{lemma}[Fixed learner depth, uniformly over the prefactor]
\label{lem:fixed-depth-uniform}
For every fixed $L\ge2$ there is a polynomial $p_L$ such that
every width-$C$, depth-$L$ tanh teacher with weights bounded
by $C$ and global signed margin greater than $1/2$ is learned
from
\[
 n\ge \varepsilon^{-1}
 p_L\bigl(d,C,1+\log(1/\varepsilon),\log(2/\delta)\bigr)
\]
examples at every $N\ge C$ and every $0<t\le1$.
The polynomial is independent of $N,t$.
\end{lemma}
\begin{proof}
We give the target-cost bound, and state the additional intrinsic
learning input precisely. Write
$a_{N,t}= -N^{-1}\log\mu_{N,t}(A_f)$.
The cloning cylinder of \cite[Proposition \texttt{prop:cloning}]{meanfield}
extends to all $0<t\le1$ as follows.
Triple the teacher readout and partition the $N$ neurons into
$C$ groups of sizes $k_b$ differing by at most one.
The exact clone has first rows equal to the teacher rows,
internal entries $B^j_{ba}/k_a$, and readout entries $v_b/k_b$.
Constrain all scaled first-layer coordinates, the internal
row projections onto the group-constant subspace, and all
scaled readout coordinates to intervals of radius $\rho t$
about this clone; the respective scalings are
$\sqrt d\,W^1$, $\sqrt N\,W^j$, and $Nv$.
There are
$NM$, where $M=d+2+(L-2)C$, independent constrained Gaussian
coordinates of variance $t^2$, and their squared centers sum
to at most $NB_0$, where
\[
 B_0=d(d+1)C^2+2(L-2)C^4+36C^4.
\]
The free internal matrices vanish on the group-constant
subspace. An elementary sphere-net Gaussian estimate bounds
all their operator norms by $b_Lt$ with conditional
probability at least $1/2$, where
$b_L=\sqrt{16\log9+8\log(4L)}$.

Put $A=2+3C^2+b_L$ and
$\rho=[2+9C^2(d+L+1)A^{L-2}]^{-1}$.
The normalized hidden error obeys
$e_1\le\rho(d+1)$ and
$e_j\le(\sqrt2C^2+\sqrt C+b_L)e_{j-1}+\rho\sqrt C$.
These inequalities hold for every $t\le1$ because both
the cylinder errors and the free operator norms decrease
with $t$. The output error is at most
$3\sqrt2C^2e_{L-1}+\rho\le1/2$, so the strict target
margin is preserved. The scalar Gaussian interval estimate
\[
 \Pr\{|Z-u|\le\rho t\}
 \ge \frac{2\rho}{\sqrt{2\pi}}
             \exp(-u^2/t^2-\rho^2)
 \quad (Z\sim\mathcal N(0,t^2))
\]
therefore gives
\[
 t^2a_{N,t}
 \le B_0+M\log\frac{e\sqrt{2\pi}}{2\rho}+\log2
 =:Q_L(d,C).
\]
This is polynomial at each fixed $L$, uniformly in $N,t$.

The imported normalized intrinsic theorem
\cite[Proposition \texttt{prop:intrinsic}]{uniformtargets}
gives a sample bound polynomial in
$d,L,1+t^2a_{N,t}$ and
\[
 \max\{1,t^H\Gain((a_{N,t}+t^{-2})/H)^{H/2}\},
 \qquad H=L-1,
\]
apart from a factor
$\varepsilon^{-1}(1+\log(1/\varepsilon))^c$
and an additive $C\varepsilon^{-1}\log(2/\delta)$.
It holds at every realizable width. Since
$\Gain(D)\le4D+2$, the squared one-layer factor is at most
$4(Q_L+1)/H+2$. Its $H/2$ power is polynomial for fixed
$L$. This proves the claim, including the smaller widths
handled by the intrinsic theorem's finite-width VC bound.
\end{proof}

\begin{proposition}[A sufficient collapse]
\label{prop:tc0-sufficient}
If $\Ppoly=\TC^0$, then there is one polynomial sample bound
jointly in $d,C,L,\varepsilon^{-1},\log(2/\delta)$ for the
teachers in Corollary~\ref{cor:collapse}, simultaneously
at every $N\ge C$ and every $0<t\le1$.
\end{proposition}
\begin{proof}
Let $R=d+C+L+2$. Lemma~\ref{lem:finite-description} below
gives a Boolean circuit for each such teacher with at most
$R^q$ gates, where $q$ is universal. Its finite description
length is polynomial in $R$, even if the teacher's real
weights have arbitrary precision.

Consider one universal Boolean circuit-evaluation family:
given a circuit description and its input assignment,
evaluate the circuit, and output zero on malformed
descriptions. This family lies in $\Ppoly$, indeed in
uniform polynomial time. Under the assumed equality it
has one fixed majority depth $h$ and one fixed polynomial
size bound in the description length. Hardcode the
teacher circuit description and leave its $d$ Boolean
inputs free. Layering and carrying constants changes $h$
by only a fixed factor. This produces a threshold circuit
of depth at most $h$ with width and coefficient magnitudes
polynomial in $R$, and normalized gate margins at least
one. The usual integer form of majority gates has all
these properties. Crucially, the same depth and polynomial
degree work for every teacher because the assumption was
applied to one universal family.

For $L\ge h+1$, apply
\cite[Theorem \texttt{thm:threshold}]{uniformtargets},
including its short-padding case. Its resource exponent
$C(h+1)^2$ is fixed, so its bound is jointly polynomial
in the stated resources and statistics.
It holds at every realizable width; each $N\ge C$ is
realizable by padding the original teacher. In particular,
the learner width need not exceed the compiled circuit
width. For the finitely many depths $2\le L<h+1$, apply
Lemma~\ref{lem:fixed-depth-uniform} and take the largest
of their polynomial degrees and constants.
\end{proof}

Thus the two complexity assumptions bracket the desired
learning statement:
\[
 \Ppoly=\TC^0
 \ \Longrightarrow\
 \text{universal joint-polynomial learning}
 \ \Longrightarrow\
 \Ppoly=\TC^1.
\]
This does not identify either implication as an equivalence.
The sufficiency proof needs no efficient construction of the
compiled circuits or computation of the posterior.

\section{An exponential depth lower bound would separate circuit classes}
\label{sec:negative}

We first record two elementary facts needed to control the resource
and subsequence quantifiers. The first uses only exact sample
consistency, and holds for any posterior of this kind.

\begin{lemma}[Finite-domain occupancy bound]
\label{lem:occupancy}
Let $X$ be a finite domain of size $M$. A posterior supported on
classifiers agreeing with all observed labels has, for every input
law, confidence-failure mass at most $\varepsilon$ with probability
at least $1-\delta$ whenever
\[
 n\ge\frac{M\log2+\log(1/\delta)}{\varepsilon}.
\]
\end{lemma}
\begin{proof}
At every observed input the posterior correct-label probability is
one. The confidence-failure set is therefore contained in the unseen
set $U$. Any fixed subset of mass greater than $\varepsilon$ is
missed by all samples with probability at most $e^{-n\varepsilon}$.
A union bound over the $2^M$ subsets gives
$\Prb\{\mathcal D(U)>\varepsilon\}\le2^M e^{-n\varepsilon}$.
\end{proof}

\begin{lemma}[Margined tanh teachers have polynomial-bit descriptions]
\label{lem:finite-description}
Suppose a tanh teacher has depth $D$, hidden widths at most $C$,
weights bounded in magnitude by $B\ge1$, and signed output margin
greater than $1/2$ on $X_d$. Its Boolean function has a Boolean
circuit of size polynomial in
\[
 d+C+D+\log(2+B).
\]
No computability or precision restriction on the original real
weights is required.
\end{lemma}
\begin{proof}
Put $m=\max\{d+1,C,2\}$,
$b=\lceil\log_2(B+2)\rceil$, $q=\lceil\log_2m\rceil$, and
\[
 p=(D+3)(b+q+4)+10,\qquad \rho=2^{-p},\qquad A=4m(B+1).
\]
Round each weight to a multiple of $\rho$ within error $\rho$.
Evaluate each affine form exactly on dyadic hidden states, and
approximate each hidden tanh to error at most $\rho$, with its
output clipped to $[-1,1]$. If $e_j$ is the largest coordinate
error at layer $j$, the Lipschitz property gives
\[
 e_j\le mB e_{j-1}+(m+1)\rho\le A e_{j-1}+A\rho.
\]
The same bound covers the final linear layer. Since
$\log_2 A\le b+q+2$, we have $e_D\le2\rho A^D<1/16$,
preserving every output sign. The rounded parameters have
$O(b+p)$ bits, and the exact affine forms have $O(b+q+p)$
integer and fractional bits. Their denominator is $2^{2p}$.

The activation evaluator must handle large preactivations without
enumerating a range of size $B$. First inspect the sign and compare
$|u|$ to $p+10$. If $|u|\ge p+10$, saturation to the appropriate
sign has error at most $2e^{-2(p+10)}<\rho$. Otherwise use oddness
and assume $0\le u\le p+10$. Set $J=\lceil16(p+10)^2\rceil$
and evaluate
\[
 E_J=\sum_{k=0}^J\frac{(2u)^k}{k!}.
\]
The exponential-series remainder is at most
\[
 e^{2u}\left(\frac{2eu}{J+1}\right)^{J+1}<2^{-p-4}.
\]
Compute $(E_J-1)/(E_J+1)$, and round to the nearest $\rho$-grid
point, clipped to $[-1,1]$. Since $E_J\ge1$, the derivative of
this quotient in $E_J$ is at most $1/2$; the total activation error
is therefore less than $\rho$.

The Taylor terms have common denominator $2^{2pJ}J!$. Its bit
length and those of the numerators are $O(p^3+p^2\log p)$.
Ordinary integer arithmetic thus evaluates the series, division,
and rounding with polynomially many Boolean gates. There are only
polynomially many affine operations and activation calls in the
teacher resources. Since $p=O(D(\log(m+1)+\log(B+2)))$, the
claimed circuit-size bound follows. The approximation algorithm
is uniform on its rounded parameter strings, even though those
strings may be chosen nonuniformly from the original real weights.
\end{proof}

The second imported input is a weaker common form of the previously
proved circuit-target learning bounds. We restrict its statistical
parameters to fixed constants because that is all the separation
argument needs.

\begin{proposition}[Imported circuit-target upper bound]
\label{prop:target-upper}
Fix $0<\varepsilon,\delta<1$. In either of the following regimes,
there is a constant $a\ge1$ such that the assertion below holds:
\begin{enumerate}[label=(\roman*),leftmargin=*]
 \item a fixed $\zeta>0$, all scales
       $0<\sigma^2\le N^{-(1+\zeta)}$, with $a$ allowed to depend
       on $\zeta,\varepsilon,\delta$;
 \item the critical scale $\sigma^2=N^{-1}$, with $a$ allowed to
       depend on $\varepsilon,\delta$.
\end{enumerate}
If a target has a bounded-fan-in Boolean circuit of wire size $s$
and depth $h\ge1$, then at any learner depth satisfying
\[
 L\ge a(h+1)\log(d+s+L+2),
\]
a sufficient sample count for the original posterior is
\[
 n\le(d+s+L+2)^{a(h+1)}.
\]
This guarantee holds for every input law and every
$N\ge C_L(f)$, where $C_L(f)$ is the least width admitting a
depth-$L$ tanh representation with weights bounded by that width
and signed output greater than $1/2$ everywhere. In regime (i) it
is simultaneous over the displayed finite scales. In regime (ii)
it holds, more generally, at every realizable width.
\end{proposition}

This is a direct weakening of
\cite[``Precision-independent threshold-circuit learning'']{unrestricted}
in regime (i), and
\cite[``Real threshold teachers without a weight bound'']{critical}
in regime (ii). To match their conventions, layer a bounded-fan-in
circuit by inserting copy wires; its size and width grow by at most
a polynomial factor, and its depth stays $O(h)$.
The fixed statistical factors and this polynomial change are
absorbed by $a$. The source in regime (i) in fact needs only
$L\ge2h$ before this harmless weakening. A realizable width is one
at which a depth-$L$ tanh network realizes the target signs; scaling
the readout then gives strict margin $1/2$ and positive Gaussian
target mass. In particular, every width at least that of a specified
depth-$L$ teacher is realizable. The bounded-weight admissibility
condition $N\ge C_L(f)$ can be stronger than realizability; we keep
the two notions distinct. No new Gaussian target-mass claim is needed
for this proposition.

\begin{theorem}[Exponential depth hardness implies $\NC\ne\Ppoly$]
\label{thm:exp-barrier}
Fix either regime of Proposition~\ref{prop:target-upper}, and fixed
$0<\varepsilon,\delta<1$. Suppose there is a sequence of targets
$f_j:X_{d_j}\to\{-1,1\}$ with tanh teachers at the same learner
depths $L_j\to\infty$, widths at most $C_j$, weights bounded by
$C_j$, and global signed margins greater than $1/2$. Assume
\[
 d_j+C_j\le L_j^b
\]
for some fixed $b$, and that obtaining the distribution-free,
all-admissible-width confidence guarantee, over $N\ge C_{L_j}(f_j)$,
requires at least
$\exp(cL_j)$ examples along this sequence, for some fixed $c>0$.
The same conclusion applies to the weaker width range $N\ge C_j$.
Then
\[
 \boxed{\qquad\NC\ne\Ppoly.\qquad}
\]
Indeed a subsequence of the targets extends to a Boolean family in
$\Ppoly\setminus\NC$.
\end{theorem}
\begin{proof}
Lemma~\ref{lem:occupancy}, with $M=2^{d_j}$, gives a universal
$O(2^{d_j})$ sample upper bound for fixed $\varepsilon,\delta$.
The assumed lower bound therefore forces
\[
 d_j\ge\frac c{\log2}L_j-O(1).
\]
Consequently $L_j=O(d_j)$ and $C_j\le\poly(d_j)$, while also
$d_j\le\poly(L_j)$. Select an infinite subsequence with distinct
strictly increasing $d_j$. Define a Boolean family $F_d$ to equal
$f_j$ at those dimensions and a constant at every other dimension.
Lemma~\ref{lem:finite-description} puts this family in $\Ppoly$.
Nonuniformity permits both the arbitrary subsequence and the rounded
teacher coefficients; no effective selection is being asserted.

Suppose $F\in\NC$. Then, for some fixed $k$ and fixed size exponent,
the selected targets have circuits of size $s_j=\poly(d_j)$ and
depth $h_j=O(\log^k(d_j+2))$. All of $d_j,s_j$ are polynomial in
$L_j$. Hence the depth condition in
Proposition~\ref{prop:target-upper} holds for all sufficiently large
$j$, and that proposition supplies a sample upper bound
\[
 n_j\le
 \exp\bigl(O(\log^{k+1}(L_j+2))\bigr)
 =\exp(o(L_j)).
\]
It holds at every admissible width $N\ge C_{L_j}(f_j)$, in particular
throughout $N\ge C_j$, and in regime (i) at every displayed scale. This
contradicts the exponential lower bound. Thus $F\notin\NC$.
\end{proof}

The separation implication applies to nonexplicit targets as well as
explicit ones. A counting or diagonalization construction of targets
with the stated polynomial teacher resources would still produce the
separated family. The global margin is important: it is what allows
arbitrary original real coefficients to be replaced by finite
descriptions.

\begin{corollary}[Every fixed stretched-exponential rate has the same barrier]
\label{cor:stretched-exponential-barrier}
Under the teacher, statistical, and width hypotheses of
Theorem~\ref{thm:exp-barrier}, replace the sample lower bound
$\exp(cL_j)$ by $\exp(cL_j^\alpha)$, for any fixed $c,\alpha>0$.
Then the same conclusion $\NC\ne\Ppoly$ holds. In particular this
applies at every fixed Gaussian exponent $\gamma\ge1$.
\end{corollary}
\begin{proof}
The occupancy bound now forces
$d_j\ge(c/\log2)L_j^\alpha-O(1)$. Hence
$L_j=O(d_j^{1/\alpha})$, and both depth and teacher width are
polynomial in $d_j$. The finite-description and subsequence argument
still give a family in $\Ppoly$. If it belonged to some $\NC^k$,
the imported learning theorem would give
$\exp(O(\log^{k+1}(L_j+2)))=\exp(o(L_j^\alpha))$ examples,
contradicting the assumed lower bound.
\end{proof}

Thus a fixed stretched-exponential rate is already stronger than
merely disproving a polynomial sample bound. The conditional
contrapositive of Corollary~\ref{cor:collapse}, under
$\Ppoly\ne\TC^1$, supplies failure of a universal polynomial
guarantee; it supplies no stretched-exponential rate. Such a rate
would also entail the stronger separation in
Corollary~\ref{cor:stretched-exponential-barrier}.

\begin{corollary}[The lower-bound barrier is uniform over all prefactors]
\label{cor:all-prefactor-negative}
The conclusion of Corollary~\ref{cor:stretched-exponential-barrier}
also holds for an arbitrary sequence of actual prefactors
$0<t_j\le1$ and realizable learner widths, with
$\sigma_j^2=t_j^2/N_j$. No separation from $t=1$ is required.
\end{corollary}
\begin{proof}
Here we use the companion uniform positive theorem
\cite[Theorem \texttt{thm:threshold}, long-padding case]{uniformtargets}
in place of Proposition~\ref{prop:target-upper}. Its precise
specialization needed here is: for fixed statistical parameters,
a bounded-fan-in Boolean circuit of size $s$ and depth $h$ has
sample bound
\[
 (d+s+L+2)^{C(h+1)}
\]
at every realizable width and every $0<t\le1$, whenever
$L\ge C(h+1)\log(d+s+L+2)$. Boolean gate coefficients and
margins are absolute constants; layering changes only the
universal polynomial factors.
The occupancy and finite-description parts of the preceding
proof do not depend on the scale. If the resulting family were
in $\NC^k$, this uniform theorem would give
$\exp(O(\log^{k+1}L_j))$ examples at the very same actual
$N_j,t_j$, contradicting $\exp(cL_j^\alpha)$.
\end{proof}

\begin{corollary}[Polynomial coefficient bit length suffices for the barrier]
\label{cor:bitlength-negative}
Fix $c,\alpha,b>0$ and fixed accuracy and confidence parameters
in $(0,1)$. Suppose targets at learner depths $L_j\to\infty$
have tanh teachers at those same depths, with widths $C_j$ and weight magnitudes at most
$B_j\ge1$, and signed output margins greater than $1/2$, with
\[
 d_j+C_j+\log(2+B_j)\le L_j^b.
\]
If the original posterior requires $\exp(cL_j^\alpha)$ examples
for the fixed distribution-free confidence guarantee at some
actual realizable widths $N_j$ and prefactors $0<t_j\le1$, then
$\NC\ne\Ppoly$. In particular the same conclusion applies to a
lower bound on a guarantee required at every realizable width.
\end{corollary}
\begin{proof}
The occupancy bound gives $d_j=\Omega(L_j^\alpha)$, so the teacher
width, depth, and coefficient bit length are polynomial in $d_j$.
Lemma~\ref{lem:finite-description} therefore places a subsequence
of these targets in a family in $\Ppoly$. If that family belonged
to some fixed $\NC^k$, its targets would have Boolean circuits of
size polynomial in $d_j$ and depth $O(\log^k d_j)$. Since also
$d_j\le L_j^b$, the uniform target theorem used in
Corollary~\ref{cor:all-prefactor-negative} applies at the same
learner depth, width, and scale and gives
$\exp(O(\log^{k+1}L_j))=\exp(o(L_j^\alpha))$ samples, a
contradiction. The all-realizable-width version of that upper
theorem needs no condition $N_j\ge B_j$.
\end{proof}

Thus polynomially many relevant coefficient bits, rather than
polynomial coefficient magnitude, suffice for the lower-bound
barrier. This does not cover an activation that itself encodes
hard truth tables: the proof uses the polynomial-time numerical
evaluator for exact tanh. Large but polynomial-bit coefficients
do not substitute for such an activation-level lookup mechanism.

\begin{corollary}[A quantitative threshold-depth obstruction]
\label{cor:quantitative-threshold-depth}
Under the hypotheses of Corollary~\ref{cor:bitlength-negative},
suppose $0<\alpha\le1$. For every fixed $a>0$, there is a constant
$c_a>0$ such that any real threshold circuit of width at most
$L_j^a$ computing $f_j$ has depth at least
\[
 c_a\frac{L_j^\alpha}{\log(L_j+2)}
\]
for all sufficiently large $j$. Its original coefficients may be
arbitrary real numbers. The constant may depend on the fixed
resource exponents and statistical parameters.
\end{corollary}
\begin{proof}
It suffices to consider threshold depth $h_j\le L_j$. The
weight-independent threshold compiler and the uniform target
upper bound give sufficient sample count
\[
 \exp(C(h_j+1)\log R_j),\qquad
 R_j=d_j+L_j^a+h_j+L_j+2,
\]
provided $L_j\ge C(h_j+1)\log R_j$; universal polynomial
compiler overhead changes only $C$. Write
$\log R_j\le B_a\log(L_j+2)$. If
\[
 h_j+1\le\kappa L_j^\alpha/\log(L_j+2),\qquad
 \kappa\le\frac{\min\{1,c\}}{2CB_a},
\]
then the padding condition holds because $L_j^\alpha\le L_j$,
and the sample upper is at most $\exp(cL_j^\alpha/2)$.
This contradicts the assumed lower bound at the same width and
scale. Absorb the additive one in $h_j+1$ by decreasing the
constant, since $L_j^\alpha/\log(L_j+2)\to\infty$.
\end{proof}

\section{Scope and remaining questions}

The product structural theorem is unconditional and applies to arbitrary
finite depths, widths, and neuronwise zero-preserving $1$-Lipschitz
activations. Its statistical consequence uses the exact original tanh
posterior at finite Gaussian scale. The product route covers every
fixed gap $\gamma\ge1+\zeta$ without a resource dependence on
the learner depth in the circuit conclusion.
Canonical features, near-baseline action, and the block cover
give a complementary converse uniform over all $0<t\le1$,
equivalently all $0<\sigma^2\le1/N$, with the depth padding,
depth-dependent circuit resources, and finite-width threshold
in Corollary~\ref{cor:uniform-converse}.
This includes arbitrary prescribed schedules approaching the
critical endpoint and all fixed exponents $\gamma\ge1$.
The actual width is chosen before the input law, and no limiting
posterior is substituted for the original one.
Corollary~\ref{cor:all-prefactor-negative}, using the companion
uniform positive theorem, also makes the stretched-exponential
lower-bound barrier uniform throughout this scale range.

The two complexity consequences address different strengths of
learning claim. A universal joint polynomial bound would imply
$\Ppoly=\TC^1$, while the stronger assumption
$\Ppoly=\TC^0$ suffices for it. Even a fixed stretched-exponential lower bound
$\exp(\Omega(L^\alpha))$, $\alpha>0$, with polynomial teacher
resources would imply $\Ppoly\ne\NC$. Neither implication
is an equivalence. Failure of a universal polynomial bound need not
give an exponential lower bound, and a universal subexponential
bound need not give the polynomial circuit sizes required for
$\TC^1$ or $\NC^2$. Familiar targets already in $\NC$, including
bounded-depth formulas and polylogarithmic-depth circuit families,
cannot establish the stated exponential lower bound because
Proposition~\ref{prop:target-upper} gives subexponential sample bounds
for them.

Several representation questions remain meaningful between these
barriers. Which subclasses of $\TC^1$ have polynomial-product tanh
representations? Can a useful quantitative relation be proved between
parallel-circuit resources and the smallest product representation?
Can one obtain a general learning bound between polynomial and
exponential depth dependence without forcing polynomial-size
parallel circuits? The present results delimit these questions but
do not settle them.

\appendix
\section{Constant-width bounded-weight teachers can encode every small circuit}
\label{app:constant-width}

The teacher width in Corollary~\ref{cor:collapse} need not grow.
The following finite-horizon construction makes this precise. The idea of
encoding discrete memory in real registers is classical: saturated-linear
neural simulations use base encodings~\cite{siegelmann-sontag}, and
Kilian and Siegelmann~\cite{kilian-siegelmann} establish sigmoidal
universality. We give an independent finite feedforward construction with
explicit depth and error bounds. In particular, we do not infer a
polynomial-time tanh simulation from an earlier universality theorem.

\begin{theorem}[A constant-width tanh compiler]
\label{thm:constant-width-compiler}
Let $d\ge1$, and suppose an acyclic Boolean circuit with $S$ gates,
each of fan-in at most two, computes $f:X_d\to\{-1,1\}$.
There is a tanh network with augmented input $(1,x)$, no separate
hidden biases, hidden width $24$, at most
\[
 H\le5000(d+S+1)^3
\]
hidden layers, every physical coefficient of absolute value at most one,
and scalar linear readout satisfying $f(x)z(x)>1$ for every $x\in X_d$.
All coefficients may be chosen dyadic with
$O(H+\log(d+1))$ bits each.
The construction only guarantees a Frobenius product bounded by
$\exp(O((d+S+1)^3))$; it is not a polynomial-product representation.
\end{theorem}

\subsection{Registers and stable amplification}

Put $n=d+S$, $M=n+1$, and set
\[
 c=\tanh1,\quad \kappa=c^{-1},\quad
 \delta=\eta=4^{-20M},\quad
 \alpha=\frac{\operatorname{arctanh}\delta}{\delta},\quad K=500M.
\]
We first use six logical registers $R,T,U,B,C,D$.
The last two are constant rails, initialized from the augmented input
by preactivations $1$ and $\operatorname{arctanh}\delta$.
The updates
\[
 C\leftarrow\tanh(\kappa C),\qquad
 D\leftarrow\tanh(\alpha D)
\]
keep them exactly equal to $c$ and $\delta$.
Here $0.761<c<0.762$ and $\kappa,\alpha<2$.
The $D$ rail is included to make tiny constants explicit; the particular
memory updates below can also use tiny coefficients multiplying $C$.

Let $\varphi(u)=\tanh(\kappa u)$.
Its fixed points include $0,\pm c$. Concavity gives
\[
 \frac{\varphi(u)}u\ge
 \frac{\tanh(3/4)}{3c/4}>1.1
 \quad(0<u\le3c/4),
\]
and on the invariant interval $[3c/4,1]$,
\[
 |\varphi'(u)|\le\frac{\operatorname{sech}^2(3/4)}c<0.8.
\]
By oddness the same bounds hold with signs reversed. Thus, whenever
$\delta/10\le|u|\le1$, $K$ iterations give
\[
 |\varphi^{K}(u)-c\sign u|\le\eta.
\]
Indeed it suffices to take
\[
 \left\lceil\frac{\log(8/\delta)}{\log1.1}\right\rceil+
 \left\lceil\frac{\log(1/\eta)}{\log1.25}\right\rceil<500M
\]
iterations. Once a register is within $\eta$ of $\pm c$, applying
$\varphi$ keeps it there. We call this a robust bit register.
Small memory registers instead use the hold update $u\leftarrow\tanh u$.
Every layer updates all registers; the schedule below specifies which
registers use a small hold, a robust hold, a reset, or an arithmetic update.

\subsection{Base-four memory and a fixed read schedule}

Store bits $b_1,\ldots,b_m\in\{-1,1\}$, with $m\le n$, as
\[
 r^*=\delta\sum_{i=1}^m b_i4^{-i}.
\]
Then $|r^*|\le\delta/3$ and, for a nonempty word,
$|r^*|\ge\delta/6$, with sign equal to $b_1$.
The ideal operations are
\[
 \text{pop}:\quad r^*\mapsto4r^*-\delta b_1,
 \qquad
 \text{insert}:\quad r^*\mapsto(r^*+\delta b)/4.
\]
Initialize
\[
 R=\tanh\left(\delta\sum_{i=1}^d x_i4^{-i}\right),
 \qquad T=U=B=0.
\]
The initial memory error is at most $\delta^3$.

To read position $p$ of $R$, first set $T\leftarrow\tanh R$.
For $i=1,\ldots,p$, set $U\leftarrow\tanh T$ and apply $K$
amplifier steps to $U$. If $i<p$, pop using
\[
 T\leftarrow\tanh\bigl(4T-(\delta/c)U\bigr).
\]
Throughout the read, hold $R$ in its small mode. The error bound below
shows that each amplification reads the intended bit, so at completion
$U$ is within $\eta$ of $c$ times the requested bit. Circuit addresses
and gate order are known in advance: there is no input-dependent control
flow. A read uses at most $1+n(K+2)$ layers.

For two operand bits, read the first and save it by
$B\leftarrow\varphi(U)$. Hold $B$ robustly while reading the second
into $U$. A NAND operation has update
\[
 U\leftarrow\tanh(C-B-U),
\]
followed by $K$ amplifier layers. Its ideal preactivation values are
$-c,c,3c$; the input error is at most $2\eta$. Hence its sign is the
NAND bit with a fixed margin, and amplification restores $\eta$ accuracy.
Insert its output by
\[
 R\leftarrow\tanh\bigl((R+(\delta/c)U)/4\bigr).
\]
Only outputs of the original circuit gates are inserted, so the word
never exceeds $d+S=n$ digits.

For completeness, an arbitrary two-input gate needs at most five NAND
operations on $B,U,T$. During this finite program $T$ is a robust bit
register; the next read overwrites it. In the following table,
$\mathsf N(X,Y)$ denotes NAND followed by $K$ amplification steps;
unwritten occupied robust registers are held by $\varphi$.
The initial operands are $B=a,U=b$, and the output is in $U$.
Truth strings are ordered $00,01,10,11$, where $0$ represents $-1$.
\begin{center}\small
\begin{tabular}{c l}
\hline
Truth string & Sequential assignments\\\hline
0000 & $T=\mathsf N(B,B);\ U=\mathsf N(B,T);\ U=\mathsf N(U,U)$\\
0001 & $U=\mathsf N(B,U);\ U=\mathsf N(U,U)$\\
0010 & $U=\mathsf N(U,U);\ U=\mathsf N(B,U);\ U=\mathsf N(U,U)$\\
0011 & $U=\mathsf N(B,B);\ U=\mathsf N(U,U)$\\
0100 & $B=\mathsf N(B,B);\ U=\mathsf N(B,U);\ U=\mathsf N(U,U)$\\
0101 & no assignment\\
0110 & $T=\mathsf N(B,U);\ B=\mathsf N(B,T);\ U=\mathsf N(U,T);\ U=\mathsf N(B,U)$\\
0111 & $B=\mathsf N(B,B);\ U=\mathsf N(U,U);\ U=\mathsf N(B,U)$\\
1000 & $B=\mathsf N(B,B);\ U=\mathsf N(U,U);\ U=\mathsf N(B,U);\ U=\mathsf N(U,U)$\\
1001 & $T=\mathsf N(B,U);\ B=\mathsf N(B,T);\ U=\mathsf N(U,T);\ U=\mathsf N(B,U);\ U=\mathsf N(U,U)$\\
1010 & $U=\mathsf N(U,U)$\\
1011 & $B=\mathsf N(B,B);\ U=\mathsf N(B,U)$\\
1100 & $U=\mathsf N(B,B)$\\
1101 & $U=\mathsf N(U,U);\ U=\mathsf N(B,U)$\\
1110 & $U=\mathsf N(B,U)$\\
1111 & $T=\mathsf N(B,B);\ U=\mathsf N(B,T)$\\\hline
\end{tabular}
\end{center}
A repeated source uses coefficient $-2$, within the bound four on
logical matrix entries. Literal constants can replace a read by
$U\leftarrow\varphi(\pm C)$; no extra memory digit is needed.
Two reads, a save, five NAND operations, and insertion cost at most
$(2n+5)K+4n+9$ layers per original gate. Including initialization and
an optional final read of an original input, the total is at most
\[
 (2n^2+6n)K+4n^2+11n+2<5000(n+1)^3=:T_{\max}.
\]
Gates outside the designated output's ancestry may first be deleted,
and the remaining gates topologically ordered with the output last.
Constant circuits and an output equal to an original input are covered
by the direct constant assignment and optional read respectively.

\subsection{Uniform error bounds}

For every real $z$,
\[
 |\tanh z-z|\le |z|^3/3,
\]
because $z-\tanh z$ is the integral of $\tanh^2$ and
$|\tanh z|\le|z|$. Whenever a small register is within $\delta/100$
of its ideal value, its magnitude is below $\delta/2$.
A hold or copy then introduces at most $\delta^3$ error.
A valid insertion satisfies
\[
 e_R^{\mathrm{new}}\le e_R/4+\delta\eta/2+\delta^3.
\]
Dropping the beneficial factor $1/4$, the main-memory error is at most
\[
 E_R=(T_{\max}+1)\delta^3+n\delta\eta.
\]
For a temporary read, a pop satisfies
\[
 e_T^{\mathrm{new}}\le4e_T+2\delta\eta+\delta^3.
\]
There are at most $n$ pops, and $T_{\max}$ bounds all intervening holds.
Consequently
\[
 E_T=4^n\bigl[(2T_{\max}+2)\delta^3+3n\delta\eta\bigr]
\]
bounds every temporary-memory error. With $M=n+1$,
\[
 \frac{E_T}{\delta}
 \le(10000M^3+2)4^{-39M}+3M4^{-19M}<\frac1{100}.
\]
For example, $M\le2^M$, $M^3\le8^M$, and $M\ge2$ give the last
inequality with ample slack. Also $E_R\le E_T$.

These estimates close a simultaneous induction over the fixed schedule.
At extraction, the ideal memory has magnitude at least $\delta/6$,
its error is below $\delta/100$, and the initialization
$U\leftarrow\tanh T$ introduces at most $\delta^3$ error.
Thus its sign is correct and its magnitude exceeds $\delta/10$.
The amplifier estimate applies, which in turn justifies the pop and
insertion recurrences. The robust NAND computation has the separate
constant-scale sign margin established above. Resetting $T$ at each
read prevents its intervening robust mode from carrying small-memory
error into that read. Every estimate is uniform on the entire Boolean
cube, without a probability or genericity assumption.

\subsection{Physical width, coefficient size, and finite precision}

Every logical layer is a linear map on six registers followed by tanh,
with matrix entries bounded by four. Replace each logical register by
four identical neurons. Replace an entry $A_{ij}$ by $A_{ij}/4$ on
every source-clone/destination-clone edge. Each destination clone then
has exactly the logical preactivation, and every physical coefficient
is bounded by one. The four initial clones use the initial row directly;
all its coefficients are already at most one. This gives width $24$
without any separate hidden bias.

The final $U$ is within $\eta$ of $cf(x)$. Read out $2U/c$, using
coefficient $1/(2c)<1$ on each of its four clones. The signed output is
at least $2(1-\eta/c)>1.9$.
Uniform cloning preserves each internal Frobenius norm: the cloned
matrix is $A\otimes(J_4/4)$ and $\norm{J_4/4}_F=1$.
The first norm is at most three, the readout norm below two, and each
internal norm at most $24$. Thus the product is at most
$6\,24^{T_{\max}}$, with no uncharged norm reduction.

Exact transcendental rails are unnecessary for the conclusion.
Round every physical coefficient toward zero to dyadic error at most
$\xi$, preserving its bound one. With $H$ hidden layers, maximum
coordinate errors obey
\[
 e_1\le(d+1)\xi,\qquad e_j\le24e_{j-1}+24\xi.
\]
The resulting readout error is at most
$(d+3)24^{H+1}\xi$. Taking
\[
 \xi\le\bigl[20(d+3)24^{H+1}\bigr]^{-1}
\]
changes every output by at most $1/20$. The signed margin remains
greater than one, and $O(H+\log(d+1))$ bits per coefficient suffice.
This proves Theorem~\ref{thm:constant-width-compiler}.

\begin{corollary}[Fixed-width teachers and the learning barrier]
\label{cor:constant-width-collapse}
Boolean families represented with global unit margin by width-$24$ tanh
networks of polynomial depth and coefficient magnitudes at most one
are exactly $\Ppoly$.
Consequently, the joint-polynomial-learning hypothesis of
Corollary~\ref{cor:collapse} already implies $\Ppoly=\TC^1$
when restricted to these teachers, provided the learner guarantee still
holds at every width $N\ge24$.
\end{corollary}
\begin{proof}
The compiler gives the inclusion of $\Ppoly$. The reverse inclusion
is Lemma~\ref{lem:finite-description}, with constant width and weight
bound. The compiler has polynomial depth and can be padded to at least
d hidden layers by robust holds before readout. Its fixed width and
coefficient bounds remain unchanged. The hypothesized sample count is
then polynomial in d, so the proof of Corollary~\ref{cor:collapse}
applies at a sufficiently large learner width chosen before the input law.
\end{proof}

\begin{remark}[The roles of precision and learner width]
The memory uses exponentially small state differences and exponentially
fine absolute weight precision, but only polynomially many description
bits. It is sensitive to polynomial-scale perturbations and gives no
polynomial-product or robust implementation of arbitrary small circuits.
At fixed learner width $24$, ordinary finite-architecture capacity bounds
already give polynomial sample bounds. The collapse conclusion depends
on a guarantee uniform over arbitrarily large learner widths, even
though the teacher width is fixed.
\end{remark}

\section{Exact proof dependencies and verification scope}
\label{app:dependencies}

The structural product theorem, canonical-feature construction, block-cover
argument, finite-event extraction, two-majority reductions, and the preceding
constant-width compiler are proved in this manuscript. The statistical
applications additionally use target-mass and intrinsic-learning results.
The following map specifies those inputs; it is not a claim that a short
companion statement replaces its supporting proof.

\begin{enumerate}[leftmargin=*]
\item The above-one product ceiling is Theorem
\texttt{thm:posterior-energy-ceiling} of
\path{unrestricted-primary-posterior-ceiling.tex} in the repository's
\path{works-in-progress/unrestricted-depth/source/} directory.
The statement, including its actual finite-scale width threshold, is
reproduced as Theorem~\ref{thm:ceiling} here.
\item The scalar Bellman baseline, small-state bridge, monotonicity, and
one-step dissipation are Definition \texttt{def:critical-scalar-value},
Lemma \texttt{lem:scalar-value-properties}, and equation
\texttt{eq:scalar-step-dissipation} in
\path{works-in-progress/critical-depth/source/critical-rate-foundations.tex}.
Proposition~\ref{prop:scalar-facts} gives the precise imported statements.
\item The finite-width Fourier target mass is Corollary
\texttt{cor:critical-finite-fourier-mass} in
\path{critical-finite-width-channel.tex}, in that same critical-depth
source directory. Its seed is constructed in
\path{critical-controlled-fourier-seeds.tex}; the finite channel proof
uses bounded-score concentration, conditional Gaussian weight shifts,
and binary relative entropy. These are finite-width inputs, not a
large-deviation limit substituted for the posterior.
\item The uniform target-mass bound is Proposition \texttt{prop:long}
and the Fourier-seed section of the companion
\path{uniform-target-classes.tex}. Its scale-change proof is detailed in
\path{nearcritical-target-reuse.md}, Sections 1--6: Gaussian reference
change; residual monotonicity; zero-cost score downgrade; the circuit
prefix; stable lifting of the scalar path; and finite-width normalized
slack. Its critical program inputs are the preceding critical-depth
sources. The short-padding proof appears directly in the companion's
Section 4.
\item The normalized intrinsic input used for sufficiency is Proposition
\texttt{prop:intrinsic} of that companion. Its two proofs are
\path{uniform-scale-intrinsic.md} and
\path{uniform-sharp-intrinsic.md}. They combine
\path{finite-width-dual-certificate.md}, the scaled importance compression
in \path{fixed-t-compression.md}, and the scalar gain estimate in
\path{uniform-baseline-gain.md}. The sharper argument uses vanishing
readout repair so that it retains the coefficient one on $V_*$.
The finite-witness class is defined by representability on every witness
of a prescribed size; it is not assumed that every function outside an
arbitrary VC class has a small separating witness.
\item The older bounded-fan-in target upper bounds used in
Proposition~\ref{prop:target-upper} are identified by these source-label pairs:
\begin{itemize}
\item \path{unrestricted-primary-threshold-circuits.tex}:\\
\path{cor:threshold-circuit-learning};
\item \path{critical-all-width-target-learning.tex}:\\
\path{cor:critical-all-width-real-threshold-learning}.
\end{itemize}
They occur in the respective source directories already specified. The all-prefactor extension instead uses
Theorem \texttt{thm:threshold} of the uniform companion.
\end{enumerate}

The named companion proof notes should accompany the manuscript whenever
its all-prefactor statistical claims are distributed. The repository
bibliography below identifies the older public sources; the proof notes
provide the additional scale-uniform arguments.
The geometric and compression Lean modules associated with this research
verify components of the structural argument. They do not constitute a
Lean proof of the full posterior converse, its target-mass inputs, the
constant-width compiler, or the complexity-class implications.
All conditional implications and the open status of unrestricted-depth
learning are unchanged by this distinction.

\begin{thebibliography}{99}
\bibitem{unrestricted}
ChatGPT-6 Astra, Kaarel H\"anni, and Dmitry Vaintrob.
\emph{Representation constraints and polynomial-depth learning in
Gaussian Bayesian networks}.
Research draft, 19 October 2026.
Repository: \url{https://github.com/kaarelh/AI-works}, directory
\path{neural-network-generalization/works-in-progress/unrestricted-depth/source/}.
The imported posterior ceiling is
\path{unrestricted-primary-posterior-ceiling.tex}, theorem
\texttt{thm:posterior-energy-ceiling}, including its explicit width
threshold. The imported target-learning bound is
\path{unrestricted-primary-threshold-circuits.tex}, corollary
\texttt{cor:threshold-circuit-learning}.

\bibitem{critical}
ChatGPT-6 Astra, Kaarel H\"anni, and Dmitry Vaintrob.
\emph{Depth and learning bounds for critical Bayesian neural networks}.
Research draft, 19 October 2026.
Repository: \url{https://github.com/kaarelh/AI-works}, directory
\path{neural-network-generalization/works-in-progress/critical-depth/source/}.
The imported scalar facts are in
\path{critical-rate-foundations.tex}, definition
\texttt{def:critical-scalar-value}, lemma
\texttt{lem:scalar-value-properties}, and equation
\texttt{eq:scalar-step-dissipation}. The imported finite-width
Fourier target mass is in
\path{critical-finite-width-channel.tex}, corollary
\texttt{cor:critical-finite-fourier-mass}. The imported learning bound is
\path{critical-all-width-target-learning.tex}, corollary
\texttt{cor:critical-all-width-real-threshold-learning}.

\bibitem{uniformtargets}
\emph{Uniform prior-scale bounds for structured Boolean targets}.
Companion research manuscript, 24 October 2026,
\path{uniform-target-classes.tex}.
The imported finite-width target mass is Proposition
\texttt{prop:long}, together with its section on Fourier seeds.
The physical scalar-program change of reference and normalized
finite-width comparison are detailed in the accompanying research
note \path{nearcritical-target-reuse.md}.

\bibitem{meanfield}
ChatGPT-6 Astra, Kaarel H\"anni, and Dmitry Vaintrob.
\emph{Mean-field Bayesian circuit learning at every width
above target complexity}.
Research draft, 15 October 2026.
Repository: \url{https://github.com/kaarelh/AI-works}, directory
\path{neural-network-generalization/main-results/meanfield-all-width/source/}.
The cloning cylinder is Proposition \texttt{prop:cloning}
in \path{cloning.tex}.

\bibitem{siegelmann-sontag}
Hava T. Siegelmann and Eduardo D. Sontag.
\emph{On the Computational Power of Neural Nets}.
Journal of Computer and System Sciences 50(1):132--150, 1995.
\url{https://doi.org/10.1006/jcss.1995.1013}.

\bibitem{kilian-siegelmann}
Joe Kilian and Hava T. Siegelmann.
\emph{The Dynamic Universality of Sigmoidal Neural Networks}.
Information and Computation 128(1):48--56, 1996.
\url{https://doi.org/10.1006/inco.1996.0062}.

\bibitem{grs}
Noah Golowich, Alexander Rakhlin, and Ohad Shamir.
\emph{Size-Independent Sample Complexity of Neural Networks}.
Conference on Learning Theory, PMLR 75, 2018.
\url{https://proceedings.mlr.press/v75/golowich18a.html}.
Full version: \url{https://arxiv.org/abs/1712.06541}.

\bibitem{sellke}
Mark Sellke.
\emph{On Size-Independent Sample Complexity of ReLU Networks}.
2024 revision.
\url{https://arxiv.org/abs/2306.01992}.

\bibitem{vardi-shamir}
Gal Vardi and Ohad Shamir.
\emph{Neural Networks with Small Weights and Depth-Separation Barriers}.
Advances in Neural Information Processing Systems, 2020.
\url{https://arxiv.org/abs/2006.00625}.

\bibitem{vardi-benign}
Gal Vardi, Daniel Reichman, Toniann Pitassi, and Ohad Shamir.
\emph{Size and Depth Separation in Approximating Benign Functions
with Neural Networks}.
Conference on Learning Theory, PMLR 134:4195--4223, 2021.
\url{https://proceedings.mlr.press/v134/vardi21a.html}.

\bibitem{jacot-rank}
Arthur Jacot.
\emph{Implicit Bias of Large Depth Networks: a Notion of Rank
for Nonlinear Functions}.
International Conference on Learning Representations, 2023.
\url{https://arxiv.org/abs/2209.15055}.

\bibitem{jacot-bottleneck}
Arthur Jacot.
\emph{Bottleneck Structure in Learned Features:
Low-Dimension vs Regularity Tradeoff}.
Advances in Neural Information Processing Systems, 2023;
2024 revision. \url{https://arxiv.org/abs/2305.19008}.
\end{thebibliography}
\end{document}
